% !TEX root = ../../main.tex
\setchapterimage{figure/chapter4封面图.jpg}
\setchapterintro{}
\chapter{线性空间与线性映射}
\section{线性空间与向量}

\begin{definition}[线性空间]
	取一个域 $\F$，$V$ 是一个非空集合。若在 $V$ 上有两种运算：加法和数乘，分别满足
	\[
	+ : V\times V\longrightarrow V,
	\qquad
	\cdot : \F\times V\longrightarrow V,
	\quad (k,\alpha)\longmapsto k\alpha,
	\]
	并且满足如下八条性质：
	\begin{enumerate}
		\item $\forall \alpha,\beta\in V$，有
		\[
		\alpha+\beta=\beta+\alpha .
		\]
		\item $\forall \alpha,\beta,\gamma\in V$，有
		\[
		(\alpha+\beta)+\gamma=\alpha+(\beta+\gamma) .
		\]
		\item $\exists 0\in V$，使得 $\forall \alpha\in V$，有
		\[
		0+\alpha=\alpha .
		\]
		\item $\forall \alpha\in V$，$\exists \beta\in V$，使得
		\[
		\alpha+\beta=0 .
		\]
		此时记 $\beta=-\alpha$。
		若 $\alpha+\beta_1=0$ 且 $\alpha+\beta_2=0$，则
		\[
		\beta_1=\beta_1+0
		=\beta_1+\alpha+\beta_2
		=0+\beta_2=\beta_2,
		\]
		故负元唯一。
		\item $\forall \alpha\in V$，有
		\[
		1\alpha=\alpha .
		\]
		\item $\forall k,l\in\F$，$\forall \alpha\in V$，有
		\[
		k(l\alpha)=(kl)\alpha .
		\]
		\item $\forall \alpha,\beta\in V$，$\forall k\in\F$，有
		\[
		k(\alpha+\beta)=k\alpha+k\beta .
		\]
		\item $\forall k,l\in\F$，$\forall \alpha\in V$，有
		\[
		(k+l)\alpha=k\alpha+l\alpha .
		\]
	\end{enumerate}
	则称 $V$ 是 $\F$ 上的一个线性空间，或称三元组 $(V,+,\cdot)$ 是一个 $\F$-线性空间。
\end{definition}

结合我们已经学过的内容，得到了如下例子。
\begin{example}
	\begin{enumerate}
		\item $\F[x]$ 与 $\F[x_1,\cdots,x_n]$ 关于多项式的加法与数乘是 $\F$ 上的线性空间。
		\item $C[a,b]$ 关于函数的加法与数乘是 $\mathbb{R}$ 上的线性空间。
		\item $\F^n$ 关于向量的加法与数乘是 $\F$ 上的线性空间。
		\item $M_{m\times n}(\F)$ 关于矩阵的加法与数乘是 $\F$ 上的线性空间。
		\item $\Hom(\F^n,\F^m)$ 关于映射的加法与数乘是 $\F$ 上的线性空间。
		\item 单元素集 $\{0\}$：定义
		\[
		0+0=0,
		\qquad
		\forall k\in\F,\quad k0=0.
		\]
		在这种情况下显然满足上述八条，仍是 $\F$-线性空间，称为零空间，记为 $0$。
		\item 对 $\mathbb{R}^{+}$，若考虑通常的加法与数乘，则 $\mathbb{R}^{+}$ 不是线性空间。
		若定义
		\[
		a\oplus b=ab,
		\qquad
		k\circ a=a^k,
		\]
		则 $(\mathbb{R}^{+},\oplus,\circ)$ 是一个实线性空间。
		\item 对 $\F\subset \mathbb{K}$，$\mathbb{K}$ 上关于数的加法与数乘是 $\F$ 上的线性空间。
		\item 对 $\F\subset \mathbb{K}$，若 $V$ 是 $\mathbb{K}$ 上的线性空间，则 $V$ 也是 $\F$ 上的线性空间。但是反之不成立。
	\end{enumerate}
\end{example}


下面，我们研究线性空间中向量与向量组的基本性质。其中一部分性质因为与第二章中所讨论的内容一致，证明也相近，故这些证明都略去，可详见第二章。

\begin{lemma}[数乘消去律]
	$V$ 是 $\F$ 上的线性空间。若 $\alpha\in V$，$k\in\F$，且
	\[
	k\alpha=0,
	\]
	则 $k=0$ 或 $\alpha=0$。
\end{lemma}
\begin{corollary}
	设 $V$ 是 $\F$ 上的线性空间。若 $\F$ 为无限域，则 $V$ 要么是单点集，要么是无限集。
\end{corollary}


与在 $\F^n$ 中的情形一致，我们定义线性组合、线性表示与线性相关，这里不再赘述。

\begin{proposition}
	设 $V$ 是 $\F$ 上的线性空间，$\beta,\alpha_1,\cdots,\alpha_s\in V$。下列三者之一成立：
	\begin{enumerate}
		\item $\beta$ 不能由 $\alpha_1,\cdots,\alpha_s$ 表出；
		\item $\beta$ 能由 $\alpha_1,\cdots,\alpha_s$ 唯一表出；
		\item $\beta$ 能由 $\alpha_1,\cdots,\alpha_s$ 无穷地表出。
	\end{enumerate}
\end{proposition}

特别地，$V$ 中零向量可以由任意向量组线性表示。

\begin{corollary}
	设 $\alpha_1,\cdots,\alpha_s\in V$，则关于 $0$ 由 $\alpha_1,\cdots,\alpha_s$ 线性表示有如下两种情形之一：
	\begin{enumerate}
		\item 若
		\[
		\sum_{i=1}^{s}k_i\alpha_i=0,
		\]
		则
		\[
		k_i=0,
		\qquad i=1,\cdots,s;
		\]
		\item $0$ 能由 $\alpha_1,\cdots,\alpha_s$ 无穷地表出。
	\end{enumerate}
\end{corollary}

下面给出一些命题，其内容与第二章完全相同，同样的方法证明。

\begin{proposition}
	设 $V$ 是 $\F$ 上的线性空间。
	\begin{enumerate}
		\item $\alpha_1,\cdots,\alpha_s\in V$ 线性相关，当且仅当至少存在一个 $\alpha_i$ 可以由其余向量线性表出。
		\item $\alpha_1,\cdots,\alpha_s\in V$ 线性无关，且 $\alpha_1,\cdots,\alpha_s,\beta$ 线性相关，则 $\beta$ 可由 $\alpha_1,\cdots,\alpha_s$ 表出。
		\item 设 $V$ 中有两个向量$ (\mathrm{I})\quad \alpha_1,\cdots,\alpha_r,
		\qquad
		(\mathrm{II})\quad \beta_1,\cdots,\beta_s.$
	
		若每个 $\alpha_i$ 都可以由 $(\mathrm{II})$ 表出，且 $r>s$，则 $(\mathrm{I})$ 线性相关。
		\item 设 $V$ 中有两个向量组
		$ (\mathrm{I})\quad \alpha_1,\cdots,\alpha_r,
		\qquad
		(\mathrm{II})\quad \beta_1,\cdots,\beta_s.$
		若每个 $\alpha_i$ 都可以由 $(\mathrm{II})$ 表出，且 $(\mathrm{I})$ 线性无关，则
		\[
		r\leq s.
		\]
		\item 设 $\alpha_1,\cdots,\alpha_m\in V$ 线性无关，$A=(a_{ij})\in M_{m\times s}(\F)$。定义
		\[
		\beta_i=a_{1i}\alpha_1+\cdots+a_{mi}\alpha_m,
		\qquad i=1,\cdots,s.
		\]
		也即
		\[
		(\beta_1\ \cdots\ \beta_s)
		=(\alpha_1\ \cdots\ \alpha_m)A_{m\times s}.
		\]
		则
		\[
		r(\beta_1,\cdots,\beta_s)=r(A).
		\]
	\end{enumerate}
\end{proposition}

同样的，我们可以定义出 $V$ 中向量组的极大无关组与秩。

接下来，给出一些之前没有的概念。 
\begin{definition}
	设 $V$ 是 $\F$ 上的线性空间，$S\subset V$。若 $S$ 中任意有限多个向量都线性无关，则称 $S$ 是线性无关集；否则称 $S$ 是线性相关集。
\end{definition}

\begin{definition}
	对 $\alpha\in V$，若存在 $S$ 中有限多个向量，使得 $\alpha$ 可由这些向量线性表示，则称 $\alpha$ 可由 $S$ 表出。也即存在一族系数 $(y_\beta)_{\beta\in S}$，其中只有有限个 $y_\beta\neq0$，使得
	\[
	\alpha=\sum_{\beta\in S}y_\beta\beta.
	\]
\end{definition}

\begin{proposition}
	设 $\alpha\in V$ 可以由 $S\subset V$ 表出，则表示方式唯一，当且仅当 $S$ 线性无关。
\end{proposition}

\begin{proof}
	对
	\[
	\alpha=\sum_{\beta\in S}y_\beta\beta,
	\]
	只有有限个 $y_\beta$ 不为 $0$。
	
	充分性：设
	\[
	\alpha=\sum_{i=1}^{r}x_i\gamma_i
	=\sum_{j=1}^{s}y_j\delta_j,
	\qquad
	\gamma_i,\delta_j\in S.
	\]
	那么
	\[
	\sum_{i=1}^{r}x_i\gamma_i-
	\sum_{j=1}^{s}y_j\delta_j=0.
	\]
	将 $\gamma_1,\cdots,\gamma_r,\delta_1,\cdots,\delta_s$ 中相同的向量合并。由 $S$ 线性无关可知，每个向量前的系数都为 $0$，故两种表示相同。
	
	必要性：若 $S$ 线性相关，则存在互不相同的 $\gamma_1,\cdots,\gamma_r\in S$ 与不全为 $0$ 的 $x_1,\cdots,x_r\in\F$，使得
	\[
	\sum_{i=1}^{r}x_i\gamma_i=0.
	\]
	因此
	\[
	\alpha=
	\sum_{\beta\in S}y_\beta\beta
	+
	\sum_{i=1}^{r}x_i\gamma_i.
	\]
	右侧给出了与原表示不同的另一种表示，与表示唯一矛盾。所以 $S$ 线性无关。
\end{proof}

\section{基与维数}

设 $V$ 是 $\F$ 上的线性空间。

\begin{definition}
	若 $V$ 的一个子集 $S$ 满足：
	\begin{enumerate}
		\item $S$ 线性无关；
		\item $\forall \alpha\in V$，$\alpha$ 可由 $S$ 表出，
	\end{enumerate}
	则称 $S$ 是 $V$ 的一个基，其中 $S$ 中元素称为基向量。零空间没有基。
\end{definition}

\begin{example}
	\begin{enumerate}
		\item $\F^n$ 中 $\varepsilon_1,\cdots,\varepsilon_n$ 构成 $\F^n$ 的一个基。
		\item $M_{m\times n}(\F)$ 中所有的 $E_{ij}\in M_{m\times n}(\F)$ 构成 $M_{m\times n}(\F)$ 的一个基。
		\item $1,x,\cdots,x^n,\cdots$ 构成 $\F[x]$ 的一个基。
	\end{enumerate}
\end{example}

\begin{proposition}(基的存在性定理)
	任意非零线性空间都有基。
\end{proposition}
证明涉及 Zorn 引理。这里只是给出一个泛泛的证明，相关语言不必完全理解，只需知道它是用 Zorn 引理证明的即可。
\begin{proof}
	设 $V$ 中所有线性无关子集构成集合
	\[
	\mathcal{F}=\{A\subset V: A\text{ 线性无关}\}.
	\]
	由于 $V$ 非零，存在 $\alpha\neq 0$，$\alpha\in V$，于是 $\{\alpha\}$ 线性无关，故 $\mathcal{F}$ 非空。
	
	在 $\mathcal{F}$ 上定义偏序关系为集合的包含关系，即
	\[
	A\leq B\Longleftrightarrow A\subset B.
	\]
	任取 $\mathcal{F}$ 中的一个链 $\mathcal{C}$，令
	\[
	M=\bigcup_{A\in\mathcal{C}}A.
	\]
	只需证明 $M\in\mathcal{F}$。任取 $M$ 中任意有限个向量 $v_1,\cdots,v_n$，设它们分别属于
	\[
	A_1,\cdots,A_n\in\mathcal{C}.
	\]
	由于 $\mathcal{C}$ 是链，在 $A_1,\cdots,A_n$ 中存在一个最大的元素，记为 $A_k$，则所有 $v_i\in A_k$。又 $A_k\in\mathcal{F}$，所以 $v_1,\cdots,v_n$ 线性无关，从而 $M\in\mathcal{F}$。
	
	因此 $(\mathcal{F},\subset)$ 的每条链都有上界。由 Zorn 引理，$\mathcal{F}$ 存在极大元 $B$。下面证明 $B$ 是 $V$ 的基。因为 $B\in\mathcal{F}$，所以 $B$ 线性无关。只需证明 $B$ 生成 $V$。
	
	反证，若存在 $v\in V$ 不能由 $B$ 表出，则 $v\notin B$，且 $B\cup\{v\}$ 仍线性无关，从而 $B\cup\{v\}\in\mathcal{F}$，这与 $B$ 的极大性矛盾。故 $B$ 生成 $V$，于是 $B$ 是 $V$ 的一个基。
\end{proof}

下面介绍基的扩充定理。证明基的存在性定理与基的扩充定理等价，需要用到商空间，只给出必要定义，建议在完成商空间的学习后再回看。

\begin{definition}
	$V$ 是 $\F$ 上的线性空间，$W$ 是 $V$ 的一个子空间。定义 $V$ 上的等价关系
	\[
	u\sim v\Longleftrightarrow u-v\in W.
	\]
	对 $v\in V$，定义
	\[
	v+W=\{v+w:w\in W\}.
	\]
	这正是上述 $\sim$ 的等价类，称为关于 $W$ 的陪集。
	
	所有这些陪集构成的集合
	\[
	V/W=\{v+W:v\in V\}
	\]
	称为商空间。
\end{definition}

\begin{proposition}(基的扩充定理)
	$V$ 是非零线性空间，$S$ 是 $V$ 中线性无关集，则 $S$ 可以扩充为 $V$ 的一个基。
	也即，存在 $C\subset V$，使得 $S\cup C$ 为 $V$ 的一个基。
\end{proposition}

\begin{proof}
	只需证扩充性与存在性的等价即可。
	\begin{enumerate}
		\item 扩充性 $\Rightarrow$ 存在性。
		
		空集 $\varnothing$ 是线性无关的，由扩充性定理，$\varnothing$ 可以扩充成为 $V$ 的一个基，记为 $B$，存在性得证。
		
		\item 存在性 $\Rightarrow$ 扩充性。
		
		考虑商空间
		\[
		V/\langle S\rangle,
		\]
		其中 $\langle S\rangle$ 表示由 $S$ 张成的子空间。由存在性定理，$V/\langle S\rangle$ 存在一个基
		\[
		\{\overline{v_i}:i\in I\},
		\qquad
		\overline{v_i}=v_i+\langle S\rangle.
		\]
		下证
		\[
		S\cup\{v_i:i\in I\}
		\]
		是 $V$ 的一个基。
		
		首先证明张成性。对任意 $v\in V$，由于 $\{\overline{v_i}:i\in I\}$ 是商空间的基，存在有限多个 $a_i\in\F$，使得
		\[
		\overline{v}=\sum_i a_i\overline{v_i}.
		\]
		于是
		\[
		v+\langle S\rangle
		=\sum_i a_i\bigl(v_i+\langle S\rangle\bigr),
		\]
		故
		\[
		v-\sum_i a_i v_i\in\langle S\rangle.
		\]
		于是存在有限多个 $b_j\in\F$ 与 $s_j\in S$，使得
		\[
		v=\sum_i a_i v_i+\sum_j b_j s_j.
		\]
		张成性得证。
		
		再证明线性无关性。设
		\[
		\sum_i a_i v_i+\sum_j b_j s_j=0.
		\]
		移项得
		\[
		\sum_i a_i v_i\in\langle S\rangle.
		\]
		于是
		\[
		\sum_i a_i\overline{v_i}=0.
		\]
		由于 $\{\overline{v_i}:i\in I\}$ 线性无关，故 $a_i=0$。那么
		\[
		\sum_j b_j s_j=0.
		\]
		又 $S$ 线性无关，故 $b_j=0$。因此 $S\cup\{v_i:i\in I\}$ 线性无关。
	\end{enumerate}
\end{proof}

\begin{lemma}
	设 $V$ 有一个基含 $n$ 个向量，那么 $V$ 的任意基都含 $n$ 个向量。
\end{lemma}

\begin{proof}
	设 $\alpha_1,\cdots,\alpha_n$ 构成 $V$ 的一个基。设 $S$ 是 $V$ 的一个基，下证 $|S|=n$。
	
	若 $|S|\geq n+1$，则 $S$ 中有 $n+1$ 个向量。由基本引理，这 $n+1$ 个向量线性相关，与 $S$ 线性无关矛盾，故 $|S|\leq n$。
	
	若 $|S|<n$，由于 $S$ 是 $V$ 的一个基，$\alpha_1,\cdots,\alpha_n$ 均可由 $S$ 表出。又 $n>|S|$，由基本引理，$\alpha_1,\cdots,\alpha_n$ 线性相关，矛盾。故 $|S|=n$。
\end{proof}

\begin{definition}
	若 $V$ 有 $n$ 个向量构成的基，则称 $V$ 是 $\F$ 上的 $n$ 维线性空间，$n$ 是 $V$ 在 $\F$ 上的维数，记为
	\[
	\dim_{\F}V=n
	\qquad (\dim V=n).
	\]
	否则称 $V$ 是 $\F$ 上的无穷维线性空间，记为
	\[
	\dim_{\F}V=\infty.
	\]
	规定 $\dim 0=0$。
\end{definition}

对 $\dim V<\infty$，$V$ 的基不写成集合的形式，而写成 $\alpha_1,\cdots,\alpha_n$ 的形式。基依赖于向量的排序顺序，例如 $\alpha_1,\alpha_2$ 与 $\alpha_2,\alpha_1$ 是不同的基。

对无限维线性空间，我们有如下命题。
\begin{proposition}
	\[
	\dim V=\infty
	\Longleftrightarrow
	\forall k\in\mathbb{Z}_{+},\ V\text{ 中有 }k\text{ 个线性无关的向量}.
	\]
\end{proposition}

\begin{example}
	\begin{enumerate}
		\item $\dim_{\F}\F^n=n$。
		\item $\dim M_{m\times n}(\F)=mn$。
		\item $W$ 是 $AX=0$ 的解空间，则 $\dim W=n-r(A)$。
		\item $\dim_{\mathbb{C}}\mathbb{C}=1$，$\dim_{\mathbb{R}}\mathbb{C}=2$。
		\item $\dim \F[x]=\infty$。
		\item $\dim C[a,b]=\infty$。
		\item $\dim_{\mathbb{R}}\mathbb{R}=1$，$\dim_{\mathbb{Q}}\mathbb{R}=\infty$。
	\end{enumerate}
\end{example}

\begin{proposition}
	\[
	\dim V\geq k
	\Longleftrightarrow
	V\text{ 中存在 }k\text{ 个线性无关的向量}.
	\]
\end{proposition}

\begin{proof}
	必要性由定义显然。
	
	充分性：若 $\dim V<k$，由基本引理，任意 $k$ 个向量线性相关，与存在 $k$ 个线性无关的向量矛盾。
\end{proof}

换一个说法，
\[
\dim V\leq k
\Longleftrightarrow
V\text{ 中任意 }k+1\text{ 个向量线性相关}.
\]

\begin{corollary}
	设 $\dim V=n$，$\alpha_1,\cdots,\alpha_n\in V$，则
	\[
	\alpha_1,\cdots,\alpha_n\text{ 是 }V\text{ 的基}
	\Longleftrightarrow
	\alpha_1,\cdots,\alpha_n\text{ 线性无关}.
	\]
\end{corollary}

\begin{proof}
	必要性由定义显然。
	
	充分性：存在 $\beta_1,\cdots,\beta_n$，它们是 $V$ 的一个基。由于 $\alpha_1,\cdots,\alpha_n$ 线性无关，且 $\dim V=n$，故$\alpha_1,\cdots,\alpha_n$与$\beta_1,\cdots,\beta_n$等价,从而它们也是 $V$ 的一个基。
\end{proof}

\begin{corollary}
	设 $\dim V=n$，$\alpha_1,\cdots,\alpha_n$ 是 $V$ 的一个基。设 $A\in M_n(\F)$ 是 $n$ 阶方阵，令
	\[
	(\beta_1\ \cdots\ \beta_n)=(\alpha_1\ \cdots\ \alpha_n)A.
	\]
	那么 $\beta_1,\cdots,\beta_n$ 是 $V$ 的基，当且仅当 $A$ 可逆。
\end{corollary}

\begin{proof}
	必要性：若 $\beta_1,\cdots,\beta_n$ 是 $V$ 的基，则 $\beta_1,\cdots,\beta_n$ 线性无关，所以
	\[
	r(\beta_1,\cdots,\beta_n)=r(A)=n.
	\]
	故 $A$ 可逆。
	
	充分性：若 $A$ 可逆，则 $r(A)=n$，于是
	\[
	r(\beta_1,\cdots,\beta_n)=n,
	\]
	故 $\beta_1,\cdots,\beta_n$ 线性无关，从而是 $V$ 的一个基。
\end{proof}

\begin{proposition}
	设 $\alpha_1,\cdots,\alpha_n$ 是 $V$ 的一个基。则对任意 $\alpha\in V$，存在唯一的 $x\in\F^n$，
	\[
	x=\begin{pmatrix}
		x_1\\
		\vdots\\
		x_n
	\end{pmatrix},
	\]
	使得
	\[
	\alpha=\sum_{i=1}^{n}x_i\alpha_i.
	\]
	即
	\[
	V=\left\{\sum_{i=1}^{n}x_i\alpha_i:x_i\in\F,\ 1\leq i\leq n\right\}.
	\]
\end{proposition}

\begin{proof}
	一方面，
	\[
	V\subset\left\{\sum_{i=1}^{n}x_i\alpha_i:x_i\in\F\right\},
	\]
	由基的定义得证。
	
	另一方面，
	\[
	\left\{\sum_{i=1}^{n}x_i\alpha_i:x_i\in\F\right\}\subset V,
	\]
	由 $V$ 是线性空间得证。
\end{proof}

\begin{definition}
	上述向量 $x\in\F^n$ 定义为 $\alpha\in V$ 在基 $\alpha_1,\cdots,\alpha_n$ 下的坐标，$x_i$ 称为第 $i$ 个坐标分量。
	
	这样，取定 $\alpha_1,\cdots,\alpha_n$ 后，可得一个一一映射
	\[
	\rho:V\longrightarrow\F^n,
	\qquad
	\alpha\longmapsto \alpha\text{ 在基 }\alpha_1,\cdots,\alpha_n\text{ 下的坐标}.
	\]
	称 $\rho$ 为关于 $\alpha_1,\cdots,\alpha_n$ 的坐标映射。
\end{definition}

\begin{lemma}
	设 $\alpha_1,\cdots,\alpha_n$ 与 $\beta_1,\cdots,\beta_n$ 是 $V$ 的两个基，则存在唯一的 $A\in\GL_n(\F)$，使得
	\[
	(\beta_1\ \cdots\ \beta_n)=(\alpha_1\ \cdots\ \alpha_n)A.
	\]
	称 $A$ 为从 $\alpha_1,\cdots,\alpha_n$ 到 $\beta_1,\cdots,\beta_n$ 的过渡矩阵。
\end{lemma}

\begin{proof}
	存在性：设 $\beta_i$ 在基 $\alpha_1,\cdots,\alpha_n$ 下的坐标为
	\[
	x_i=\begin{pmatrix}
		a_{1i}\\
		\vdots\\
		a_{ni}
	\end{pmatrix},
	\qquad i=1,\cdots,n.
	\]
	令 $A=(x_1\ \cdots\ x_n)=(a_{ji})$，则
	\[
	(\beta_1\ \cdots\ \beta_n)=(\alpha_1\ \cdots\ \alpha_n)A.
	\]
	由于 $\beta_1,\cdots,\beta_n$ 是 $V$ 的基，故
	\[
	r(\beta_1,\cdots,\beta_n)=r(A)=n,
	\]
	所以 $A$ 可逆，即 $A\in\GL_n(\F)$。
	
	唯一性：若
	\[
	(\alpha_1\ \cdots\ \alpha_n)A=(\alpha_1\ \cdots\ \alpha_n)B,
	\]
	则对任意 $i$，有
	\[
	\sum_{j=1}^{n}a_{ji}\alpha_j
	=
	\sum_{j=1}^{n}b_{ji}\alpha_j.
	\]
	于是
	\[
	\sum_{j=1}^{n}(a_{ji}-b_{ji})\alpha_j=0.
	\]
	由 $\alpha_1,\cdots,\alpha_n$ 线性无关，得 $a_{ji}=b_{ji}$，故 $A=B$。
\end{proof}

这样，取定 $\alpha_1,\cdots,\alpha_n$ 后，又可定义如下映射：
\[
\sigma:\{V\text{ 的基}\}\longrightarrow \GL_n(\F),
\]
\[
\beta_1,\cdots,\beta_n
\longmapsto
\text{从 }\alpha_1,\cdots,\alpha_n\text{ 到 }\beta_1,\cdots,\beta_n\text{ 的过渡矩阵 }A.
\]

因此有如下推论。
\begin{corollary}
	设 $\alpha_1,\cdots,\alpha_n$ 与 $\beta_1,\cdots,\beta_n$ 是 $V$ 的两个基。设
	\[
	\alpha=(\alpha_1\ \cdots\ \alpha_n)x
	=(\beta_1\ \cdots\ \beta_n)y
	=(\alpha_1\ \cdots\ \alpha_n)Ay,
	\]
	其中 $A$ 为从 $\alpha_1,\cdots,\alpha_n$ 到 $\beta_1,\cdots,\beta_n$ 的过渡矩阵，则
	\[
	x=Ay.
	\]
\end{corollary}

\section{子空间}

\begin{definition}[子空间]
	设 $V$ 是 $\F$ 上的线性空间，$W$ 是 $V$ 的非空子集。若 $W$ 对 $V$ 运算封闭，即
	\[
	\forall u,v\in W,\quad u+v\in W,
	\qquad
	\forall k\in\F,\ \forall u\in W,\quad ku\in W,
	\]
	则称 $W$ 是 $V$ 的一个子空间。
	

\end{definition}
	子空间 $W$ 有以下性质：
\begin{enumerate}
	\item $0\in W$；
	\item 若 $U$ 是 $W$ 的子空间，$W$ 是 $V$ 的子空间，则 $U$ 是 $V$ 的子空间；
	\item 若 $W$ 是 $V$ 的子空间，则将 $V$ 上的运算限制在 $W$ 上，使得 $W$ 成为一个 $\F$-线性空间。
\end{enumerate}
\begin{example}
	\begin{enumerate}
		\item 令
		\[
		\F[x]_n=\{f\in\F[x]:f^{(n)}(x)=0\}.
		\]
		则
		\[
		\dim\F[x]_n=n,
		\]
		其基为
		\[
		1,x,\cdots,x^{n-1}.
		\]
		并且
		\[
		\F[x]_n\subset \F[x]_{n+1}\subset\cdots,
		\]
		$\F[x]_n$ 是 $\F[x]_{n+1}$ 的子空间。
		
		\item 设
		\[
		f\in\Hom(\F^n,\F^m).
		\]
		定义
		\[
		\ker f=\{\alpha\in\F^n:f(\alpha)=0\},
		\]
		\[
		\Im(f)=\{\beta\in\F^m:\exists \alpha\in\F^n,\ f(\alpha)=\beta\}.
		\]
		则 $\ker f$ 是 $\F^n$ 的子空间，$\Im(f)$ 是 $\F^m$ 的子空间。
		
		\item $\mathbb{R}^{+}$ 不是 $\mathbb{R}$ 的子空间。
	\end{enumerate}
\end{example}

\begin{definition}[平凡子空间]
	设 $V$ 是线性空间，则 $\{0\}$ 与 $V$ 本身都是 $V$ 的子空间，称为 $V$ 的平凡子空间。
\end{definition}

\begin{lemma}
	若 $W$ 是 $V$ 的子空间，则
	\[
	\dim W\leq \dim V.
	\]
\end{lemma}

\begin{proof}
	$W$ 中任意线性无关向量组在 $V$ 中仍线性无关，因此 $W$ 中线性无关向量组的长度不超过 $V$ 中线性无关向量组长度的上界，故
	\[
	\dim W\leq \dim V.
	\]
\end{proof}

注意：
\begin{enumerate}
	\item 若 $\alpha_1,\cdots,\alpha_n$ 是 $V$ 的一个基，$W$ 是 $V$ 的一个 $m$ 维子空间，则 $\alpha_1,\cdots,\alpha_n$ 中的 $m$ 个向量不一定可以构成 $W$ 的一个基。
	
	\item 若 $W$ 是 $V$ 的一个真子空间，则不一定有
	\[
	\dim W<\dim V.
	\]
	例如
	\[
	W=\left\{\sum_{i=1}^{s}a_i x^i:s\in\mathbb{Z}_{+},\ a_i\in\F\right\},
	\qquad
	V=\F[x].
	\]
	则 $W\subsetneq V$，但
	\[
	\dim W=\dim V=\infty.
	\]
\end{enumerate}

\begin{proposition}
	设 $\dim V<\infty$，$U$ 是 $V$ 的子空间。若
	\[
	\dim U=\dim V,
	\]
	则
	\[
	U=V.
	\]
	其等价形式为：若 $U$ 是 $V$ 的真子空间，则
	\[
	\dim U<\dim V.
	\]
\end{proposition}

\begin{proof}
	设 $U$ 是 $V$ 的真子空间。令
	\[
	\alpha_1,\cdots,\alpha_n
	\]
	为 $U$ 的一个基。由于 $U\neq V$，存在 $\sigma\in V\setminus U$，那么 $\sigma$ 不能由 $\alpha_1,\cdots,\alpha_n$ 表出，故
	\[
	\alpha_1,\cdots,\alpha_n,\sigma
	\]
	线性无关。因此 $V$ 中至少有 $n+1$ 个线性无关向量，于是
	\[
	\dim V>n=\dim U.
	\]
	这与 $\dim U=\dim V$ 矛盾，故 $U=V$。
\end{proof}

\begin{definition}[余维数]
	$W$ 为 $V$ 的子空间。称
	\[
	\dim V-\dim W
	\]
	为 $W$ 在 $V$ 中的余维数，记为
	\[
	\operatorname{codim}_{V}W.
	\]
	这是商空间 $V/W$ 的维数。
\end{definition}

下面考虑向量组生成的子空间。

\begin{definition}
	设 $V$ 是 $\F$ 上的线性空间，$\alpha_1,\cdots,\alpha_s\in V$，令
	\[
	W=\left\{\alpha\in V:\exists k_i\in\F,\ \alpha=\sum_{i=1}^{s}k_i\alpha_i\right\}.
	\]
	显然 $W$ 也是 $V$ 的子空间，称 $W$ 为由 $\alpha_1,\cdots,\alpha_s$ 线性生成的子空间，记为
	\[
	\langle \alpha_1,\cdots,\alpha_s\rangle
	\quad\text{或}\quad
	\Span\{\alpha_1,\cdots,\alpha_s\}.
	\]
	若 $U$ 是 $V$ 的子空间，
	\[
	U=\langle \beta_1,\cdots,\beta_t\rangle,
	\]
	则称 $\beta_1,\cdots,\beta_t$ 为 $U$ 的一组生成元。
\end{definition}

\begin{lemma}
	设 $\alpha_1,\cdots,\alpha_s\in V$。
	\begin{enumerate}
		\item
		\[
		\dim\langle \alpha_1,\cdots,\alpha_s\rangle\leq s,
		\]
		且 $\alpha_1,\cdots,\alpha_s$ 的任一个极大无关子组构成 $\langle \alpha_1,\cdots,\alpha_s\rangle$ 的一个基。
		
		\item
		\[
		\langle \alpha_1,\cdots,\alpha_s\rangle
		=
		\langle \beta_1,\cdots,\beta_t\rangle
		\]
		当且仅当 $\alpha_1,\cdots,\alpha_s$ 与 $\beta_1,\cdots,\beta_t$ 等价。
	\end{enumerate}
\end{lemma}

下面，我们用一种绕开 Zorn 引理的方法考虑有限维线性空间基的扩充定理

\begin{proposition}
	设 $\dim V<\infty$，$S\subset V$。则存在有限集 $S'\subset V$，使得
	\[
	\langle S\cup S'\rangle=V.
	\]
\end{proposition}

\begin{proof}
	令
	\[
	W=\langle S\rangle.
	\]
	若
	\[
	\dim V-\dim W=0,
	\]
	则 $W=V$，取 $S'=\varnothing$ 即可。
	
	设
	\[
	\dim V-\dim W=m>0.
	\]
	取 $W$ 的一个基$ \beta_1,\cdots,\beta_s.$
	因为 $W\neq V$，存在
	\[
	u\in V\setminus W.
	\]
	下面证明 $u,\beta_1,\cdots,\beta_s$ 线性无关。设
	\[
	x_1\beta_1+\cdots+x_s\beta_s+yu=0.
	\]
	若 $y\neq0$，则
	\[
	u=-\frac{1}{y}(x_1\beta_1+\cdots+x_s\beta_s)\in W,
	\]
	矛盾。因此 $y=0$。再因 $\beta_1,\cdots,\beta_s$ 线性无关，得
	\[
	x_1=\cdots=x_s=0.
	\]
	故 $u,\beta_1,\cdots,\beta_s$ 线性无关。
	
	设
	\[
	W'=\langle S\cup\{u\}\rangle.
	\]
	则
	\[
	\dim W'=\dim W+1.
	\]
	于是
	\[
	\dim V-\dim W'=m-1.
	\]
	对 $m$ 归纳，便可逐次加入向量 $u_1,\cdots,u_m\in V$，使得
	\[
	\langle S\cup\{u_1,\cdots,u_m\}\rangle=V.
	\]
	取
	\[
	S'=\{u_1,\cdots,u_m\},
	\]
	即得结论。
\end{proof}

\section{子空间的运算}

先考虑子空间的交、并、补等运算。

我们可以看到，两个子空间的并集一般不是子空间。

\begin{proposition}
	设 $V_1,V_2$ 是 $V$ 的子空间，则 $V_1\cup V_2$ 是 $V$ 的子空间当且仅当
	\[
	V_1\subset V_2
	\quad\text{或}\quad
	V_2\subset V_1.
	\]
\end{proposition}

\begin{proof}
	充分性显然。
	
	必要性：设 $V_1\cup V_2$ 是子空间。若
	\[
	V_1\not\subset V_2,
	\qquad
	V_2\not\subset V_1,
	\]
	则存在
	\[
	\alpha\in V_1,\quad \alpha\notin V_2,
	\qquad
	\beta\in V_2,\quad \beta\notin V_1.
	\]
	断言
	\[
	\alpha+\beta\notin V_1\cup V_2.
	\]
	否则，若 $\alpha+\beta\in V_1$，则
	\[
	\beta=(\alpha+\beta)-\alpha\in V_1,
	\]
	矛盾；若 $\alpha+\beta\in V_2$，则
	\[
	\alpha=(\alpha+\beta)-\beta\in V_2,
	\]
	矛盾。故 $\alpha+\beta\notin V_1\cup V_2$，这与 $V_1\cup V_2$ 是子空间矛盾。
	因此 $V_1\subset V_2$ 或 $V_2\subset V_1$。
\end{proof}

子空间的交集是子空间。

\begin{proposition}
	任意多个子空间的交是子空间。
\end{proposition}

\begin{proof}
	根据定义显然。
\end{proof}

\begin{definition}
	设 $V_1,\cdots,V_s$ 是 $V$ 的子空间。定义
	\[
	V_1+\cdots+V_s
	=\sum_{i=1}^{s}V_i
	=\left\{\alpha\in V:\exists \alpha_i\in V_i,\
	\alpha=\alpha_1+\cdots+\alpha_s\right\}\subset V.
	\]
	显然，$\sum\limits_{i=1}^{s}V_i$ 是 $V$ 的子空间，称 $\sum\limits_{i=1}^{s}V_i$ 是 $V_1,\cdots,V_s$ 的和。
\end{definition}

根据向量加法的交换律，子空间的加法显然有交换律，即
\[
V_1+V_2=V_2+V_1.
\]

\begin{example}
	\begin{enumerate}
		\item
		\[
		\langle \alpha_1,\cdots,\alpha_s\rangle
		+
		\langle \beta_1,\cdots,\beta_t\rangle
		=
		\langle \alpha_1,\cdots,\alpha_s,\beta_1,\cdots,\beta_t\rangle.
		\]
		
		\item 
		\[
		S(n)=\{A\in M_n(\F):A'=A\},
		\]
		\[
		A(n)=\{A\in M_n(\F):A'=-A\}.
		\]
		则
		\[
		M_n(\F)=S(n)+A(n).
		\]
	\end{enumerate}
\end{example}

注意：一般来说，
\[
V_1+V_2=V_1+V_3
\quad\not\Rightarrow\quad
V_2=V_3,
\]
即不一定有 $V_2=V_3$。例如可在 $\mathbb{R}^2$ 中取反例。
同样地，一般也有
\begin{equation}
V_1\cap (V_2+V_3)
\neq
(V_1\cap V_2)+(V_1\cap V_3).
\tag{$*$}
\end{equation}
\begin{proposition}
	设 $V_1,V_2$ 是 $V$ 的子空间，且 $\dim(V_1+V_2)<\infty$，则
	\[
	\dim(V_1+V_2)
	=
	\dim V_1+\dim V_2-\dim(V_1\cap V_2).
	\]
\end{proposition}

\begin{proof}
	设 $V_1\cap V_2\neq\{0\}$，取 $V_1\cap V_2$ 的一组基
	\[
	\varepsilon_1,\cdots,\varepsilon_r,
	\qquad
	r=\dim(V_1\cap V_2).
	\]
	由基的扩充定理，将其扩充为 $V_1$ 的一组基
	\[
	\varepsilon_1,\cdots,\varepsilon_r,
	\varepsilon_{r+1},\cdots,\varepsilon_{r+a},
	\]
	也扩充为 $V_2$ 的一组基
	\[
	\varepsilon_1,\cdots,\varepsilon_r,
	\delta_{r+1},\cdots,\delta_{r+b}.
	\]
	下证
	\[
	\varepsilon_1,\cdots,\varepsilon_r,
	\varepsilon_{r+1},\cdots,\varepsilon_{r+a},
	\delta_{r+1},\cdots,\delta_{r+b}
	\]
	是 $V_1+V_2$ 的一组基。
	
	一方面，由
	\[
	V_1+V_2
	=
	\langle \varepsilon_1,\cdots,\varepsilon_r,
	\varepsilon_{r+1},\cdots,\varepsilon_{r+a}\rangle
	+
	\langle \varepsilon_1,\cdots,\varepsilon_r,
	\delta_{r+1},\cdots,\delta_{r+b}\rangle
	\]
	可知上述向量组生成 $V_1+V_2$。
	
	另一方面，设
	\[
	\sum_{i=1}^{r}x_i\varepsilon_i
	+
	\sum_{j=1}^{a}y_j\varepsilon_{r+j}
	+
	\sum_{k=1}^{b}z_k\delta_{r+k}
	=0.
	\]
	移项得
	\[
	V_1\ni
	\sum_{i=1}^{r}x_i\varepsilon_i
	+
	\sum_{j=1}^{a}y_j\varepsilon_{r+j}
	=
	-\sum_{k=1}^{b}z_k\delta_{r+k}
	\in V_2.
	\]
	因此
	\[
	\sum_{k=1}^{b}z_k\delta_{r+k}\in V_1\cap V_2.
	\]
	故存在 $a_1,\cdots,a_r\in\F$，使得
	\[
	\sum_{k=1}^{b}z_k\delta_{r+k}
	=
	\sum_{i=1}^{r}a_i\varepsilon_i.
	\]
	由
	\[
	\varepsilon_1,\cdots,\varepsilon_r,
	\delta_{r+1},\cdots,\delta_{r+b}
	\]
	线性无关，得 $a_i=0$ 且 $z_k=0$。于是
	\[
	\sum_{i=1}^{r}x_i\varepsilon_i
	+
	\sum_{j=1}^{a}y_j\varepsilon_{r+j}=0.
	\]
	由
	\[
	\varepsilon_1,\cdots,\varepsilon_r,
	\varepsilon_{r+1},\cdots,\varepsilon_{r+a}
	\]
	线性无关，得 $x_i=y_j=0$。故该向量组线性无关，从而为 $V_1+V_2$ 的基。
	
	因此
	\[
	\dim(V_1+V_2)=r+a+b
	=(r+a)+(r+b)-r
	=\dim V_1+\dim V_2-\dim(V_1\cap V_2).
	\]
	当 $V_1\cap V_2=\{0\}$ 时，同理可证。
\end{proof}

考虑
\[
\dim(V_1+V_2+V_3)
=
\dim(V_1+V_2)+\dim V_3-
\dim((V_1+V_2)\cap V_3),
\]
由 $(*)$ 可见，对 $n\geq 3$ 的情形，$\dim(\sum_{i=1}^{n}V_i)$ 一般没有一个简单的公式。

\begin{example}
	设
	\[
	V_1=\langle \alpha_1,\cdots,\alpha_s\rangle,
	\qquad
	V_2=\langle \beta_1,\cdots,\beta_t\rangle.
	\]
	由 $\alpha_1,\cdots,\alpha_s$ 与 $\beta_1,\cdots,\beta_t$ 的极大无关组求出 $\dim V_1$ 与 $\dim V_2$，再由
	\[
	\alpha_1,\cdots,\alpha_s,
	\beta_1,\cdots,\beta_t
	\]
	的极大无关组求出 $\dim(V_1+V_2)$，则可得
	\[
	\dim(V_1\cap V_2)
	=
	\dim V_1+\dim V_2-\dim(V_1+V_2).
	\]
	进一步，若要求 $V_1\cap V_2$ 的一组基，可考虑齐次方程
	\[
	\sum_{i=1}^{s}x_i\alpha_i
	+
	\sum_{j=1}^{t}y_j\beta_j=0.
	\]
	对其每个解，都有
	\[
	\sum_{i=1}^{s}x_i\alpha_i
	=
	-\sum_{j=1}^{t}y_j\beta_j\in V_1\cap V_2,
	\]
	由这些向量中取极大无关组，即可得到 $V_1\cap V_2$ 的一组基。
\end{example}

\begin{definition}
	设 $V_1,\cdots,V_s$ 是 $V$ 的子空间。若对任意
	\[
	\alpha\in V_1+\cdots+V_s,
	\]
	存在唯一的 $\alpha_i\in V_i$，$1\leq i\leq s$，使得
	\[
	\alpha=\alpha_1+\cdots+\alpha_s,
	\]
	则称 $V_1+\cdots+V_s$ 是直和，记为
	\[
	V_1\oplus\cdots\oplus V_s.
	\]
\end{definition}

即使是直和，与交运算之间也不满足分配律。也即，若
\[
V=V_1\oplus\cdots\oplus V_s,
\]
且 $W$ 是 $V$ 的子空间，一般仍有
\[
W\cap V
=
W\cap (V_1\oplus\cdots\oplus V_s)
\neq
(W\cap V_1)\oplus\cdots\oplus(W\cap V_s).
\]

\begin{lemma}
	设 $V_1,\cdots,V_s$ 是 $V$ 的子空间。则 $\sum\limits_{i=1}^{s}V_i$ 是直和当且仅当
	\[
	0=\alpha_1+\cdots+\alpha_s,
	\quad \alpha_i\in V_i
	\quad\Longrightarrow\quad
	\alpha_1=\cdots=\alpha_s=0.
	\]
\end{lemma}

\begin{proof}
	必要性显然。
	
	充分性：对任意 $\alpha\in\sum\limits_{i=1}^{s}V_i$，设
	\[
	\alpha=\alpha_1+\cdots+\alpha_s
	=\beta_1+\cdots+\beta_s,
	\qquad
	\alpha_i,\beta_i\in V_i.
	\]
	则
	\[
	(\alpha_1-\beta_1)+\cdots+(\alpha_s-\beta_s)=0.
	\]
	由假设得
	\[
	\alpha_i-\beta_i=0,
	\qquad i=1,\cdots,s.
	\]
	故 $\alpha_i=\beta_i$，从而表示唯一。因此 $\sum\limits_{i=1}^{s}V_i$ 是直和。
\end{proof}

\begin{corollary}
	设 $V_1,V_2$ 是 $V$ 的子空间，则
	\[
	V_1+V_2\text{ 是直和}
	\quad\Longleftrightarrow\quad
	V_1\cap V_2=\{0\}.
	\]
	进一步，
	\[
	V_1+V_2\text{ 是直和}
	\quad\Longleftrightarrow\quad
	\dim(V_1\cap V_2)=0.
	\]
\end{corollary}

\begin{proof}
	充分性：设
	\[
	0=\alpha_1+\alpha_2,
	\qquad
	\alpha_i\in V_i.
	\]
	则
	\[
	\alpha_1=-\alpha_2\in V_2.
	\]
	从而
	\[
	\alpha_1\in V_1\cap V_2=\{0\},
	\]
	故 $\alpha_1=\alpha_2=0$。由引理可知，$V_1+V_2$ 是直和。
	
	必要性：若 $\beta\in V_1\cap V_2$，则
	\[
	0=\beta+(-\beta),
	\qquad
	\beta\in V_1,
	\quad
	-\beta\in V_2.
	\]
	由直和表示的唯一性得 $\beta=0$，从而 $V_1\cap V_2=\{0\}$。
\end{proof}

\begin{corollary}
	设 $V_1,\cdots,V_s$ 是 $V$ 的子空间，则 $V_1+\cdots+V_s$ 是直和当且仅当对每个 $i$ 都有
	\[
	V_i\cap\left(\sum_{j\neq i}V_j\right)=0.
	\]
\end{corollary}

\begin{proof}
	必要性：设
	\[
	\alpha\in V_i\cap\left(\sum_{j\neq i}V_j\right).
	\]
	则存在 $\alpha_j\in V_j$，使得
	\[
	\sum_{j\neq i}\alpha_j=\alpha.
	\]
	于是
	\[
	\sum_{j\neq i}\alpha_j-\alpha=0.
	\]
	由直和性及上一引理得 $\alpha=0$。
	
	充分性：设
	\[
	0=\gamma_1+\cdots+\gamma_s,
	\qquad
	\gamma_i\in V_i.
	\]
	则
	\[
	\gamma_i=-\sum_{j\neq i}\gamma_j\in\sum_{j\neq i}V_j.
	\]
	又 $\gamma_i\in V_i$，故
	\[
	\gamma_i\in V_i\cap\left(\sum_{j\neq i}V_j\right)=0,
	\]
	即 $\gamma_i=0$。由上一引理，$V_1+\cdots+V_s$ 是直和。
\end{proof}

\begin{proposition}
	设 $V_1,\cdots,V_s$ 是 $V$ 的子空间。以下三条命题等价：
	\begin{enumerate}
		\item $V_1+\cdots+V_s$ 是直和；
		\item 若 $\xi_{i1},\cdots,\xi_{it_i}$ 是 $V_i$ 的基，则
		\[
		\xi_{11},\cdots,\xi_{1t_1},
		\cdots,
		\xi_{s1},\cdots,\xi_{st_s}
		\]
		是 $V_1+\cdots+V_s$ 的基；
		\item
		\[
		\dim(V_1\oplus\cdots\oplus V_s)
		=\dim V_1+\cdots+\dim V_s.
		\]
	\end{enumerate}
\end{proposition}

\begin{proof}
	$1\Rightarrow 2$：显然
	\[
	V_1+\cdots+V_s
	=
	\langle \xi_{11},\cdots,\xi_{1t_1},
	\cdots,
	\xi_{s1},\cdots,\xi_{st_s}\rangle.
	\]
	只需证这些向量线性无关。设
	\[
	\sum_{i=1}^{s}\sum_{j=1}^{t_i}x_{ij}\xi_{ij}=0.
	\]
	令
	\[
	\alpha_i=\sum_{j=1}^{t_i}x_{ij}\xi_{ij},
	\qquad
	\alpha_i\in V_i.
	\]
	于是
	\[
	\sum_{i=1}^{s}\alpha_i=0.
	\]
	由 $V_1+\cdots+V_s$ 是直和，得
	\[
	\alpha_1=\cdots=\alpha_s=0.
	\]
	再由 $\xi_{i1},\cdots,\xi_{it_i}$ 线性无关，得
	\[
	x_{ij}=0,
	\qquad
	1\leq i\leq s,
	\quad
	1\leq j\leq t_i.
	\]
	从而无关性得证。
	
	$2\Rightarrow 1$：设
	\[
	0=\gamma_1+\cdots+\gamma_s,
	\qquad
	\gamma_i\in V_i.
	\]
	又设
	\[
	\gamma_i=\sum_{j=1}^{t_i}y_{ij}\xi_{ij},
	\qquad
	y_{ij}\in\F.
	\]
	于是
	\[
	\sum_{i=1}^{s}\sum_{j=1}^{t_i}y_{ij}\xi_{ij}=0.
	\]
	由 $2$ 中的向量组线性无关，得 $y_{ij}=0$，从而 $\gamma_i=0$。由引理可知，$V_1+\cdots+V_s$ 是直和。
	
	$2\Rightarrow 3$：若 $2$ 成立，则上述向量组是 $V_1+\cdots+V_s$ 的基，所以
	\[
	\dim(V_1\oplus\cdots\oplus V_s)
	=t_1+\cdots+t_s
	=\dim V_1+\cdots+\dim V_s.
	\]
	
	$3\Rightarrow 2$：上述向量组总能生成 $V_1+\cdots+V_s$。若其线性相关，则
	\[
	\dim(V_1+\cdots+V_s)
	=
	r\{\xi_{11},\cdots,\xi_{1t_1},\cdots,\xi_{s1},\cdots,\xi_{st_s}\}
	<t_1+\cdots+t_s,
	\]
	与 $3$ 矛盾，故该向量组线性无关，从而为一组基。
\end{proof}

\begin{definition}
	$U$ 是 $V$ 的子空间，$W$ 是 $V$ 的子空间，若
	\[
	V=U\oplus W,
	\]
	则称 $W$ 为 $U$ 在 $V$ 中的补子空间。
\end{definition}

对于非平凡子空间，补子空间一般不唯一，且可能有无穷多个。

\begin{definition}
	设 $U$ 是 $V$ 的子空间，则定义嵌入映射
	\[
	\iota:U\longrightarrow V,
	\qquad
	\alpha\longmapsto \alpha.
	\]
	称 $\iota$ 为从 $U$ 到 $V$ 的嵌入。
	
	若
	\[
	V=V_1\oplus\cdots\oplus V_s,
	\]
	则对任意 $\alpha\in V$，存在唯一的 $\alpha_i\in V_i$，使得
	\[
	\alpha=\sum_{i=1}^{s}\alpha_i.
	\]
	定义
	\[
	\pi_{V_i}:V\longrightarrow V_i,
	\qquad
	\alpha\longmapsto \alpha_i.
	\]
	称 $\pi_{V_i}$ 为从 $V$ 到 $V_i$ 的投影。
\end{definition}
\section{线性映射}
\begin{definition}
	设 $V,W$ 是 $\F$ 上的线性空间。若映射$f:V\longrightarrow W$,
	满足
	\begin{enumerate}
		\item $\forall \alpha,\beta\in V$，有
		\[
		f(\alpha+\beta)=f(\alpha)+f(\beta);
		\]
		\item $\forall \alpha\in V$，$\forall k\in\F$，有
		\[
		f(k\alpha)=k f(\alpha),
		\]
	\end{enumerate}
	则称 $f$ 是 $V$ 到 $W$ 的\textbf{线性映射}，记作
	\[
	f\in\Hom(V,W).
	\]
	若 $f\in\Hom(V,V)$，则称 $f$ 是 $V$ 上的\textbf{线性变换}，并记
	\[
	\End(V)=\Hom(V,V).
	\]
\end{definition}

规定零映射
\[
0:V\longrightarrow W,
\qquad
\alpha\longmapsto 0,
\]
是 $\Hom(V,W)$ 中的零元。

对 $f\in\Hom(V,W)$，定义
\[
\ker f=\{\alpha\in V:f(\alpha)=0\},
\]
\[
\Im(f)=\{\beta\in W:\exists\alpha\in V,\text{ 使得 }f(\alpha)=\beta\}.
\]
一般地，对 $W_1\subset W$，定义
\[
f^{-1}(W_1)=\{\alpha\in V:f(\alpha)\in W_1\},
\]
称为 $W_1$ 在 $f$ 下的原像。注意 $f^{-1}(W_1)$ 有可能等于 $\varnothing$。

对 $V_1\subset V$，定义
\[
f(V_1)=\{\beta\in W:\exists\alpha\in V_1,\text{ 使得 }\beta=f(\alpha)\},
\]
称为 $V_1$ 在 $f$ 下的像。于是
\[
\ker f=f^{-1}(0),
\qquad
\Im f=f(V).
\]

\begin{proposition}
	设 $f\in\Hom(V,W)$。则
	\begin{enumerate}
		\item $f(0)=0$；
		\item 若 $V_1$ 是 $V$ 的子空间，则 $f(V_1)$ 是 $W$ 的子空间。特别地，$\Im f$ 是 $W$ 的子空间；
		\item 若 $W_1$ 是 $W$ 的子空间，则 $f^{-1}(W_1)$ 是 $V$ 的子空间。
	\end{enumerate}
\end{proposition}

\begin{proof}
	\begin{enumerate}
		\item 由线性性，
		\[
		f(0)=f(0+0)=f(0)+f(0),
		\]
		故 $f(0)=0$。
		
		\item 设 $\alpha,\beta\in f(V_1)$，则存在 $x_1,x_2\in V_1$，使得
		\[
		\alpha=f(x_1),
		\qquad
		\beta=f(x_2).
		\]
		于是
		\[
		\alpha+\beta=f(x_1)+f(x_2)=f(x_1+x_2)\in f(V_1),
		\]
		且对任意 $k\in\F$，
		\[
		k\alpha=kf(x_1)=f(kx_1)\in f(V_1).
		\]
		因此 $f(V_1)$ 是 $W$ 的子空间。
		
		\item 设 $\alpha,\beta\in f^{-1}(W_1)$，则
		\[
		f(\alpha)\in W_1,
		\qquad
		f(\beta)\in W_1.
		\]
		由于 $W_1$ 是子空间，故
		\[
		f(\alpha+\beta)=f(\alpha)+f(\beta)\in W_1,
		\]
		从而
		\[
		\alpha+\beta\in f^{-1}(W_1).
		\]
		又对任意 $k\in\F$，
		\[
		f(k\alpha)=k f(\alpha)\in W_1,
		\]
		从而 $k\alpha\in f^{-1}(W_1)$。故 $f^{-1}(W_1)$ 是 $V$ 的子空间。
	\end{enumerate}
\end{proof}

\begin{definition}
	称 $\dim\ker f$ 为 $f$ 的\textbf{零度}，记为
	\[
	\operatorname{nullity}f=\dim\ker f.
	\]
\end{definition}

\begin{proposition}
	设 $f\in\Hom(V,W)$，则 $f$ 是单射当且仅当$ \ker f=0.$
\end{proposition}

\begin{proof}
	充分性：设 $\ker f=0$。若 $f(\alpha)=f(\beta)$，则
	\[
	f(\alpha)-f(\beta)=0,
	\]
	即
	\[
	f(\alpha-\beta)=0.
	\]
	所以 $\alpha-\beta\in\ker f=0$，从而 $\alpha=\beta$，故 $f$ 是单射。
	
	必要性：若 $f$ 是单射，且 $\alpha\in\ker f$，则
	\[
	f(\alpha)=0=f(0).
	\]
	由单射性得 $\alpha=0$，故 $\ker f=0$。
\end{proof}

\begin{corollary}
	设 $f\in\Hom(V,W)$，则 $f$ 是单射当且仅当对任意线性无关的
	$\alpha_1,\cdots,\alpha_s\in V$，都有$f(\alpha_1),\cdots,f(\alpha_s) $线性无关。
\end{corollary}
\begin{corollary}
	若存在单射 $f\in\Hom(V,W)$，且 $V$ 有限维，则
	\[
	\dim V\leq \dim W.
	\]
	反过来，当 $V,W$ 均为有限维空间时，若 $\dim V\leq \dim W$，则存在单射
	$f\in\Hom(V,W)$。
\end{corollary}

\begin{proof}
	设 $\dim V=s$，取 $V$ 的一组基$\alpha_1,\cdots,\alpha_s. $
	若 $f$ 是单射，则由上一个推论，$ f(\alpha_1),\cdots,f(\alpha_s)$
	线性无关。因此 $W$ 中有 $s$ 个线性无关向量，故
	\[
	\dim V=s\leq \dim W.
	\]
\end{proof}

\begin{corollary}
	设 $U\subset V$，则 $U$ 是 $V$ 的子空间当且仅当存在某个线性空间 $W$ 与
	$f\in\Hom(V,W)$，使得
	\[
	U=\ker f.
	\]
\end{corollary}

\begin{proof}
	充分性显然，因为 $\ker f$ 是 $V$ 的子空间。
	
	必要性：若 $U$ 是 $V$ 的子空间，取 $U$ 在 $V$ 中的一个补子空间 $W$，即
	\[
	V=U\oplus W.
	\]
	设
	\[
	\pi_W:V\longrightarrow W
	\]
	为投影映射，则
	\[
	\ker \pi_W=U.
	\]
\end{proof}

\begin{theorem}
	设 $f\in\Hom(V,W)$，且 $\dim V<\infty$。则
	\[
	\dim\ker f+\dim\Im(f)=\dim V.
	\]
\end{theorem}

\begin{proof}
	设 $\dim V=n$。取 $\ker f$ 的一组基
	$\varepsilon_1,\cdots,\varepsilon_s$，
	并将它扩充为 $V$ 的一组基
	\[
	\varepsilon_1,\cdots,\varepsilon_s,
	\varepsilon_{s+1},\cdots,\varepsilon_n.
	\]
	下面证明
	$f(\varepsilon_{s+1}),\cdots,f(\varepsilon_n)$
	是 $\Im(f)$ 的一组基。
	
	首先，
	\[
	\Im(f)
	=\left\langle f(\varepsilon_1),\cdots,f(\varepsilon_s),
	f(\varepsilon_{s+1}),\cdots,f(\varepsilon_n)\right\rangle.
	\]
	由于 $\varepsilon_1,\cdots,\varepsilon_s\in\ker f$，故
	\[
	f(\varepsilon_1)=\cdots=f(\varepsilon_s)=0.
	\]
	于是
	\[
	\Im(f)
	=\left\langle f(\varepsilon_{s+1}),\cdots,f(\varepsilon_n)\right\rangle.
	\]
	
	下面证明
	$f(\varepsilon_{s+1}),\cdots,f(\varepsilon_n)$
	线性无关。设
	\[
	a_{s+1}f(\varepsilon_{s+1})+\cdots+a_nf(\varepsilon_n)=0.
	\]
	则
	\[
	f(a_{s+1}\varepsilon_{s+1}+\cdots+a_n\varepsilon_n)=0,
	\]
	因此
	\[
	a_{s+1}\varepsilon_{s+1}+\cdots+a_n\varepsilon_n\in\ker f.
	\]
	又因为
	\[
	\ker f=\left\langle \varepsilon_1,\cdots,\varepsilon_s\right\rangle,
	\]
	所以存在标量 $b_1,\cdots,b_s$，使得
	\[
	a_{s+1}\varepsilon_{s+1}+\cdots+a_n\varepsilon_n
	=b_1\varepsilon_1+\cdots+b_s\varepsilon_s.
	\]
	移项得
	\[
	-b_1\varepsilon_1-\cdots-b_s\varepsilon_s
	+a_{s+1}\varepsilon_{s+1}+\cdots+a_n\varepsilon_n=0.
	\]
	由于
	\[
	\varepsilon_1,\cdots,\varepsilon_s,
	\varepsilon_{s+1},\cdots,\varepsilon_n
	\]
	是 $V$ 的一组基，故线性无关。因此
	\[
	a_{s+1}=\cdots=a_n=0.
	\]
	所以
	$f(\varepsilon_{s+1}),\cdots,f(\varepsilon_n)$
	线性无关。
	
	综上，
	$f(\varepsilon_{s+1}),\cdots,f(\varepsilon_n)$
	是 $\Im(f)$ 的一组基，故
	\[
	\dim\Im(f)=n-s,
	\qquad
	\dim\ker f=s.
	\]
	从而
	\[
	\dim\ker f+\dim\Im(f)=s+(n-s)=n=\dim V.
	\]
\end{proof}

\begin{corollary}
	设 $f\in\Hom(V,W)$。若
	\[
	\dim V=\dim W<\infty,
	\]
	则 $f$ 是单射当且仅当 $f$ 是满射。
\end{corollary}

\begin{proof}
	必要性：若 $f$ 单射，则
	\[
	\dim\ker f=0.
	\]
	由上一定理，
	\[
	\dim\Im f=\dim V=\dim W.
	\]
	又 $\Im f\subset W$，故 $\Im f=W$，即 $f$ 满射。
	
	充分性：若 $f$ 满射，则
	\[
	\Im f=W,
	\]
	故
	\[
	\dim\Im f=\dim W=\dim V.
	\]
	由上一定理得
	\[
	\dim\ker f=0,
	\]
	于是 $\ker f=0$，故 $f$ 单射。
\end{proof}

特别地，若 $f\in\End(V)$ 且 $\dim V<\infty$，则 $f$ 是单射当且仅当 $f$ 是满射。

\begin{proposition}
	设 $f\in\End(V)$，且 $\dim V=n$。则
	\[
	V=\ker f+\Im f
	\quad\Longleftrightarrow\quad
	\ker f+\Im f\text{ 是直和}.
	\]
\end{proposition}

\begin{proof}
	由维数公式，
	\[
	\dim(\ker f+\Im f)
	=\dim\ker f+\dim\Im f
	-\dim(\ker f\cap\Im f).
	\]
	又由上一定理，
	\[
	\dim\ker f+\dim\Im f=n=\dim V.
	\]
	因此
	\[
	\dim(\ker f+\Im f)
	=n-\dim(\ker f\cap\Im f).
	\]
	于是
	\[
	V=\ker f+\Im f
	\]
	当且仅当
	\[
	\dim(\ker f+\Im f)=n,
	\]
	当且仅当
	\[
	\dim(\ker f\cap\Im f)=0,
	\]
	当且仅当
	\[
	\ker f\cap\Im f=0.
	\]
	这正等价于 $\ker f+\Im f$ 是直和。
\end{proof}

\section{Hom空间}

\begin{definition}
	设 $V,W$ 是 $\F$ 上的线性空间，记
	\[
	\mathcal{F}(V,W)=\{f\mid f:V\longrightarrow W\}
	\]
	为从 $V$ 到 $W$ 的所有映射构成的集合。
	
	定义其上的加法与数乘如下：对 $f,g\in\mathcal{F}(V,W)$，定义
	\[
	f+g:V\longrightarrow W,
	\qquad
	\alpha\longmapsto f(\alpha)+g(\alpha),
	\]
	即
	\[
	(f+g)(\alpha)=f(\alpha)+g(\alpha).
	\]
	对 $k\in\F$，$f\in\mathcal{F}(V,W)$，定义
	\[
	kf:V\longrightarrow W,
	\qquad
	\alpha\longmapsto kf(\alpha),
	\]
	即
	\[
	(kf)(\alpha)=k f(\alpha).
	\]
	显然，$\mathcal{F}(V,W)$ 是一个线性空间，而
	$\Hom(V,W)$ 是 $\mathcal{F}(V,W)$ 的子空间。
\end{definition}

\begin{theorem}(线性映射存在性定理)
	设 $V,W$ 是 $\F$ 上的线性空间，$\dim V<\infty$。取
	$\alpha_1,\cdots,\alpha_n$ 为 $V$ 的一组基，则对任意
	$\beta_1,\cdots,\beta_n\in W$，存在唯一的
	$f\in\Hom(V,W)$，使得
	\[
	f(\alpha_i)=\beta_i,
	\qquad i=1,\cdots,n.
	\]
\end{theorem}

\begin{proof}
	存在性：对任意 $\alpha\in V$，设$ \alpha=\sum\limits_{i=1}^{n}x_i\alpha_i,$
	定义$f(\alpha)=\sum\limits_{i=1}^{n}x_i\beta_i.$
	这里 $(x_1,\cdots,x_n)'$ 为 $\alpha$ 在基
	$\alpha_1,\cdots,\alpha_n$ 下的坐标。
	
	若
	\[
	\alpha=\sum_{i=1}^{n}x_i\alpha_i,
	\qquad
	\alpha'=\sum_{i=1}^{n}z_i\alpha_i,
	\]
	则
	\[
	\alpha+\alpha'=\sum_{i=1}^{n}(x_i+z_i)\alpha_i,
	\]
	从而
	\[
	f(\alpha+\alpha')
	=
	\sum_{i=1}^{n}(x_i+z_i)\beta_i
	=
	\sum_{i=1}^{n}x_i\beta_i+
	\sum_{i=1}^{n}z_i\beta_i
	=f(\alpha)+f(\alpha').
	\]
	又对 $k\in\F$，有
	\[
	k\alpha=\sum_{i=1}^{n}kx_i\alpha_i,
	\]
	于是
	\[
	f(k\alpha)=\sum_{i=1}^{n}kx_i\beta_i
	=k\sum_{i=1}^{n}x_i\beta_i
	=kf(\alpha).
	\]
	故 $f$ 为线性映射，且显然满足
	\[
	f(\alpha_i)=\beta_i,
	\qquad i=1,\cdots,n.
	\]
	
	唯一性：设 $g\in\Hom(V,W)$ 也满足$g(\alpha_i)=\beta_i,
	\qquad i=1,\cdots,n. $则$(f-g)(\alpha_i)=0,
	\qquad i=1,\cdots,n. $
	由于 $\alpha_1,\cdots,\alpha_n$ 是 $V$ 的一组基，任意
	$\alpha\in V$ 均可写为 $\alpha=\sum\limits_{i=1}^{n}x_i\alpha_i$，于是
	\[
	(f-g)(\alpha)
	=
	\sum_{i=1}^{n}x_i(f-g)(\alpha_i)=0.
	\]
	所以 $f-g=0$，即 $f=g$。唯一性得证。
\end{proof}

\begin{corollary}
	若 $\dim V=n$，$\dim W=m$，则
	\[
	\dim \Hom(V,W)=\dim W\cdot \dim V=mn.
	\]
\end{corollary}

\begin{proof}
	取 $V$ 的一组基 $\alpha_1,\cdots,\alpha_n$，取 $W$ 的一组基
	$\beta_1,\cdots,\beta_m$。
	
	对 $1\leq i\leq m$，$1\leq j\leq n$，由上一定理，存在唯一的
	$f_{ij}\in\Hom(V,W)$，使得对 $k=1,\cdots,n$，有
	\[
	f_{ij}(\alpha_k)=
	\begin{cases}
		\beta_i, & k=j,\\
		0, & k\neq j.
	\end{cases}
	\]
	下证 $\{f_{ij}\mid 1\leq i\leq m,
	1\leq j\leq n\}$ 为 $\Hom(V,W)$ 的一组基。
	
	先证无关性。设
	\[
	\sum_{i=1}^{m}\sum_{j=1}^{n}z_{ij}f_{ij}=0.
	\]
	代入 $\alpha_k$，$1\leq k\leq n$，得
	\[
	0=
	\left(\sum_{i=1}^{m}\sum_{j=1}^{n}z_{ij}f_{ij}\right)(\alpha_k)
	=\sum_{i=1}^{m}z_{ik}\beta_i.
	\]
	由 $\beta_1,\cdots,\beta_m$ 线性无关，得
	\[
	z_{ik}=0,
	\qquad i=1,\cdots,m,
	\quad k=1,\cdots,n.
	\]
	故 $\{f_{ij}\}$ 线性无关。
	
	再证生成性。任取 $f\in\Hom(V,W)$。对每个
	$1\leq k\leq n$，存在 $y_{1k},\cdots,y_{mk}\in\F$，使得
	\[
	f(\alpha_k)=\sum_{i=1}^{m}y_{ik}\beta_i.
	\]
	断言
	\[
	f=\sum_{i=1}^{m}\sum_{j=1}^{n}y_{ij}f_{ij}.
	\]
	事实上，对 $\alpha_k$ 有
	\[
	\left(\sum_{i=1}^{m}\sum_{j=1}^{n}y_{ij}f_{ij}\right)(\alpha_k)
	=
	\sum_{i=1}^{m}y_{ik}\beta_i
	=f(\alpha_k).
	\]
	由上一定理的唯一性可知两边相等。因此 $\{f_{ij}\}$ 生成
	$\Hom(V,W)$。
	
	综上，$\{f_{ij}\}$ 是一组基，共有 $mn$ 个元素，故
	\[
	\dim\Hom(V,W)=mn.
	\]
\end{proof}

在研究映射的复合之前，我们需要知道如下事实。

\begin{proposition}\label{4.6.4}
	设有集合之间的映射
	\[
	S_1\xrightarrow{\ f\ }S_2\xrightarrow{\ g\ }S_3.
	\]
	则有：
	\begin{enumerate}
		\item 若 $gf$ 是单射，则 $f$ 是单射；
		\item 若 $gf$ 是满射，则 $g$ 是满射。
	\end{enumerate}
\end{proposition}

\begin{proof}
	1. 若 $f$ 不是单射，则存在 $y_1\neq y_2$，使得$f(y_1)=f(y_2). $于是$gf(y_1)=gf(y_2), $
	这与 $gf$ 是单射矛盾。故 $f$ 是单射。
	
	2. 若 $gf$ 是满射，则对任意 $z\in S_3$，存在 $x\in S_1$，使得$gf(x)=z.$
	令 $y=f(x)\in S_2$，则$g(y)=z.$
	故 $g$ 是满射。
\end{proof}

容易看到，对
\[
f\in\Hom(V_1,V_2),
\qquad
g\in\Hom(V_2,V_3),
\]
有
\[
gf\in\Hom(V_1,V_3).
\]
若 $f_1,f_2\in\Hom(V_1,V_2)$，$g\in\Hom(V_2,V_3)$，则
\[
g(f_1+f_2)=gf_1+gf_2.
\]
若 $f\in\Hom(V_1,V_2)$，$g_1,g_2\in\Hom(V_2,V_3)$，则
\[
(g_1+g_2)f=g_1f+g_2f.
\]
因此，线性映射的复合与加法之间满足分配律。

\begin{proposition}
	对 $\End(V)$，其上的加法、乘法（复合）以及数乘，使其成为一个结合代数。
\end{proposition}

对 $T\in\End(V)$，规定
\[
T^0=\id_V=1_V
\]
为 $V$ 上的恒等变换，并定义 $T$ 的幂：
\[
T^2=TT,
\qquad
\cdots,
\qquad
T^k=T^{k-1}T.
\]
若
\[
g(x)=\sum_{i=0}^{n}a_i x^i\in\F[x],
\]
则定义
\[
g(T)=\sum_{i=0}^{n}a_iT^i,
\]
它是一个 $V$ 上的线性变换。
\section{同构}

\begin{definition}
	对 $f\in\Hom(V,W)$，若存在
	$g\in\Hom(W,V)$，使得
	\[
	fg=\id_{W},
	\qquad
	gf=\id_{V},
	\]
	则称 $f$ 是从 $V$ 到 $W$ 的同构映射。
\end{definition}

\begin{proposition}
	对 $f\in\Hom(V,W)$，若存在
	$g\in\Hom(W,V)$，使得
	\[
	gf=\id_{V},
	\qquad
	fg=\id_{W},
	\]
	则 $g$ 唯一。
\end{proposition}

\begin{proof}
	设 $g_1,g_2\in\Hom(W,V)$，且
	\[
	g_1f=\id_{V}, fg_1=\id_{W},	g_2f=\id_{V}, fg_2=\id_{W}.
	\]
	则
	\[
	g_1
	=g_1\id_{W}
	=g_1(fg_2)
	=(g_1f)g_2
	=\id_{V}g_2
	=g_2.
	\]
	因此 $g$ 唯一。
\end{proof}

这样，$g$ 记为 $f^{-1}$，称为 $f$ 的逆映射。

\begin{definition}
	$\End(V)$ 中的同构称为 $V$ 上的自同构。例如，对
	$a\in\F$ 且 $a\neq 0$，$a\id_{V}$ 是自同构。
	
	记 $\Aut(V)$ 为 $V$ 上所有自同构的集合。可以看到，
	$\Aut(V)$ 对复合构成一个群。
\end{definition}

\begin{example}
	\begin{enumerate}
		\item $\Aut(\F)=\F\setminus\{0\}$。
		\item 设 $T\in\End(V)$，且存在 $g\in\F[x]$，使得
		\[
		g(T)=0.
		\]
		若 $h\in\F[x]$ 且 $(h,g)=1$，则
		\[
		h(T)\in\Aut(V).
		\]
		事实上，由 $(h,g)=1$，存在 $u,v\in\F[x]$，使得
		\[
		uh+vg=1.
		\]
		于是
		\[
		u(T)h(T)+v(T)g(T)=1,
		\]
		而 $g(T)=0$，故
		\[
		u(T)h(T)=1.
		\]
		从而
		\[
		u(T)=h(T)^{-1}.
		\]
	\end{enumerate}
\end{example}

\begin{lemma}
	设 $f\in\Hom(V,W)$，则 $f$ 是同构当且仅当 $f$ 单且满。
\end{lemma}

\begin{proof}
	必要性:\ref{4.6.4}
	
	充分性：若 $f$ 单且满，存在映射
	$f^{-1}:W\to V$。下面证明 $f^{-1}$ 为线性映射。
	
	任取 $\beta_1,\beta_2\in W$，设
$	f^{-1}(\beta_1)=\alpha_1,\qquad
	f^{-1}(\beta_2)=\alpha_2,$
	
	即
	\[
	f(\alpha_1)=\beta_1,\qquad
	f(\alpha_2)=\beta_2.
	\]
	于是
	\[
	f(\alpha_1+\alpha_2)=f(\alpha_1)+f(\alpha_2)=\beta_1+\beta_2,
	\]
	故
	\[
	f^{-1}(\beta_1+\beta_2)
	=
	\alpha_1+\alpha_2
	=
	f^{-1}(\beta_1)+f^{-1}(\beta_2).
	\]
	又对 $k\in\F$，有
	\[
	f(k\alpha_1)=kf(\alpha_1)=k\beta_1,
	\]
	故
	\[
	f^{-1}(k\beta_1)=kf^{-1}(\beta_1).
	\]
	因此 $f^{-1}\in\Hom(W,V)$，从而 $f$ 是同构。
\end{proof}

\begin{corollary}
	同构映射的复合也是同构映射。
\end{corollary}

\begin{definition}
	设 $V,W$ 是 $\F$ 上的线性空间。若存在同构
	$f:V\to W$，则称 $V$ 同构于 $W$，记为	$V\cong W.$
	否则称 $V$ 不同构于 $W$，记为$V\not\cong W.$
	显然，同构关系是等价关系。
\end{definition}

\begin{lemma}
	对有限维线性空间 $V$，若 $\dim V=n$，则
$V\cong \F^n.$
	
\end{lemma}

\begin{proof}
	取 $V$ 的一组基 $\alpha_1,\cdots,\alpha_n$，定义
	\[
	\varphi:V\longrightarrow \F^n,
	\qquad
	\alpha=\sum_{i=1}^{n}x_i\alpha_i
	\longmapsto
	(x_1,\cdots,x_n)'.
	\]
	不难验证 $\varphi$ 是一个同构映射。
\end{proof}

\begin{theorem}
	若 $V,W$ 为有限维线性空间，则
	\[
	V\cong W
	\quad\Longleftrightarrow\quad
	\dim V=\dim W.
	\]
\end{theorem}

\begin{proof}
	充分性：若 $\dim V=\dim W=n$
	\[V \cong \F^n \cong W\]
	必要性：设 $V\cong W$，取同构映射 $f:V\to W$。取 $V$ 的一组基
	$\alpha_1,\cdots,\alpha_n$。下面证明
	$f(\alpha_1),\cdots,f(\alpha_n)$ 是 $W$ 的一组基。
	
	任取 $\beta\in W$，由于 $f$ 满射，存在 $\alpha\in V$，使得
	$f(\alpha)=\beta$。设
	\[
	\alpha=\sum_{i=1}^{n}x_i\alpha_i,
	\]
	则
	\[
	\beta=f(\alpha)
	=f\left(\sum_{i=1}^{n}x_i\alpha_i\right)
	=\sum_{i=1}^{n}x_if(\alpha_i).
	\]
	故 $f(\alpha_1),\cdots,f(\alpha_n)$ 生成 $W$。
	

	由 $\alpha_1,\cdots,\alpha_n$ 线性无关，得
显然$f(\alpha_1),\cdots,f(\alpha_n)$ 线性无关。因此它是 $W$ 的一组基，
	从而 $\dim W=n=\dim V$。
	
	
\end{proof}

\begin{example}
	\begin{enumerate}
		\item 若 $\dim_{\F}V=n$，$\dim_{\F}W=m$，则
		\[
		\Hom(V,W)\cong \Hom(W,V)
		\cong M_{m\times n}(\F)
		\cong M_{n\times m}(\F).
		\]
		\item 记
		\[
		\F[x]_n=\{\,f\in\F[x]\mid f^{(n)}(x)=0\,\},
		\]
		则
		\[
		\F[x]_n\cong \F^n.
		\]
		\item 若 $\dim V=n$，则对偶空间
		\[
		\Hom(V,\F)\cong V.
		\]
		\item 若 $A\in M_{m\times n}(\F)$，$V$ 是方程组
		\[
		AX=0
		\]
		的解空间，则
		\[
		V\cong \F^{\,n-r(A)}.
		\]
		\item 若 $\dim V<\infty$，$W$ 是 $V$ 的子空间，且
		\[
		V=W\oplus W_1=W\oplus W_2,
		\]
		则
		\[
		W_1\cong W_2.
		\]
		\item $\mathbb{R}[x]\not\cong C[a,b]$。
		\item 若 $\dim V<\infty$，则 $V$ 不能同构于其真子空间；若
		$\dim V=\infty$，则可以取
		\[
		\F[x]\quad\text{与}\quad \F[x^2],
		\]
		显然
		\[
		\F[x]\cong \F[x^2].
		\]
	\end{enumerate}
\end{example}

\begin{proposition}
	$V$ 同构于其真子空间当且仅当
	\[
	\dim V=\infty.
	\]
\end{proposition}

\begin{proof}
	必要性：设 $V\cong W$，且 $W\subsetneqq V$。若 $\dim V<\infty$，则
	\[
	\dim W=\dim V,
	\]
	从而 $W=V$，矛盾。故
	\[
	\dim V=\infty.
	\]
	
	充分性：对无限维空间 $V$，取 $V$ 的一组基 $B$。去掉 $B$ 中的一个基向量
	$\beta$，令
	\[
	W=\langle B\setminus\{\beta\}\rangle.
	\]
	由于 $B$ 与 $B\setminus\{\beta\}$ 等势，故两者张成的空间同构。因此
	\[
	V\cong W,
	\qquad
	W\subsetneqq V.
	\]
\end{proof}

\begin{example}
	设 $A\in M_n(\F), P\in\GL_n(\F)$，且
	\[
	PAP^{-1}=
	\begin{pmatrix}
		E_r&0\\
		0&0
	\end{pmatrix}.
	\]
	记
	\[
	V=\{\,C\in M_n(\F)\mid AC=CA\,\}.
	\]
	则 $V$ 是线性空间。求 $\dim V$。
\end{example}

\begin{proof}[解]
	令
	\[
	W=\left\{
	D\in M_n(\F)\ \middle|\ 
	\begin{pmatrix}
		E_r&0\\
		0&0
	\end{pmatrix}D
	=
	D\begin{pmatrix}
		E_r&0\\
		0&0
	\end{pmatrix}
	\right\}.
	\]
	设
	\[
	D=
	\begin{pmatrix}
		D_{11}&D_{12}\\
		D_{21}&D_{22}
	\end{pmatrix},
	\]
	其中 $D_{11}\in M_r(\F)$，$D_{22}\in M_{n-r}(\F)$。由
	\[
	\begin{pmatrix}
		E_r&0\\
		0&0
	\end{pmatrix}
	D
	=
	D
	\begin{pmatrix}
		E_r&0\\
		0&0
	\end{pmatrix}
	\]
	得
	\[
	\begin{pmatrix}
		D_{11}&D_{12}\\
		0&0
	\end{pmatrix}
	=
	\begin{pmatrix}
		D_{11}&0\\
		D_{21}&0
	\end{pmatrix}.
	\]
	于是
	\[
	D_{12}=0,\qquad D_{21}=0,
	\]
	即
	\[
	D=
	\begin{pmatrix}
		D_{11}&0\\
		0&D_{22}
	\end{pmatrix}.
	\]
	因此
	\[
	\dim W=r^2+(n-r)^2=n^2-2nr+2r^2.
	\]
	
	下面证明 $W\cong V$。取
	\[
	f:V\longrightarrow W,
	\qquad
	C\longmapsto PCP^{-1}.
	\]
	对任意 $C\in V$，有
	\[
	\begin{pmatrix}
		E_r&0\\
		0&0
	\end{pmatrix}PCP^{-1}
	=
	PAP^{-1}PCP^{-1}
	=
	PACP^{-1},
	\]
	又因 $AC=CA$，得
	\[
	PACP^{-1}
	=
	PCAP^{-1}
	=
	PCP^{-1}PAP^{-1}
	=
	PCP^{-1}
	\begin{pmatrix}
		E_r&0\\
		0&0
	\end{pmatrix}.
	\]
	故 $f(C)\in W$。
	
	容易验证 $f$ 是同构映射，其逆映射为
	\[
	D\longmapsto P^{-1}DP.
	\]
	从而
	\[
	\dim V=\dim W=n^2-2nr+2r^2.
	\]
\end{proof}
\section{有限维空间之间的线性映射}

设 $\dim V=n,\dim W=m$，且 $f\in\Hom(V,W)$。
取 $V$ 的一组基 $\alpha_1,\cdots,\alpha_n$，$W$ 的一组基
$\beta_1,\cdots,\beta_m$。对每个 $i$，设
\[
f(\alpha_i)=
(\beta_1\ \cdots\ \beta_m)
\begin{pmatrix}
	a_{1i}\\
	\vdots\\
	a_{mi}
\end{pmatrix}.
\]
于是
\[
\bigl(f(\alpha_1)\ \cdots\ f(\alpha_n)\bigr)
=
(\beta_1\ \cdots\ \beta_m)
\begin{pmatrix}
	a_{11}&\cdots&a_{1n}\\
	\vdots& &\vdots\\
	a_{m1}&\cdots&a_{mn}
\end{pmatrix}
=(\beta_1\ \cdots\ \beta_m)A.
\]
其中 $A$ 的第 $i$ 列即为 $f(\alpha_i)$ 在基
$\beta_1,\cdots,\beta_m$ 下的坐标。

\begin{definition}
	称上述矩阵 $A$ 为 $f$ 在基 $\alpha_1,\cdots,\alpha_n$
	与 $\beta_1,\cdots,\beta_m$ 下的矩阵，记作
	\[
	\bigl(f(\alpha_1)\ \cdots\ f(\alpha_n)\bigr)
	=(\beta_1\ \cdots\ \beta_m)A.
	\]
	显然，也有一个映射
	\[
	\rho:\Hom(V,W)\longrightarrow M_{m\times n}(\F),
	\qquad
	f\longmapsto A.
	\]
\end{definition}

\begin{lemma}
	上述映射 $\rho$ 是一个同构映射。
\end{lemma}

\begin{proof}
	首先证明 $\rho$ 是线性映射。设 $f,g\in\Hom(V,W)$，
	且
	\[
	\rho(f)=A,
	\qquad
	\rho(g)=B.
	\]
	即
	\[
	\bigl(f(\alpha_1)\ \cdots\ f(\alpha_n)\bigr)
	=(\beta_1\ \cdots\ \beta_m)A,
	\qquad
	\bigl(g(\alpha_1)\ \cdots\ g(\alpha_n)\bigr)
	=(\beta_1\ \cdots\ \beta_m)B.
	\]
	因此
	\[
	\rho(f+g)=A+B=\rho(f)+\rho(g),
	\qquad
	\rho(kf)=k\rho(f).
	\]
	其次，若 $f\in\ker\rho$，则 $f$ 在这两组基下的矩阵为 $0$，
	从而
	\[
	f(\alpha_i)=0,
	\qquad i=1,\cdots,n.
	\]
	由于 $\alpha_1,\cdots,\alpha_n$ 是 $V$ 的一组基，故 $f=0$。
	所以 $\rho$ 是单射。
	又
	\[
	\dim\Hom(V,W)=\dim M_{m\times n}(\F)=mn,
	\]
	故 $\rho$ 也是满射，从而是同构映射。
\end{proof}

\begin{lemma}
	设 $f\in\Hom(U,V)$，$g\in\Hom(V,W)$。
	取 $U,V,W$ 的基分别为
	\[
	\alpha_1,\cdots,\alpha_l;
	\qquad
	\beta_1,\cdots,\beta_m;
	\qquad
	\gamma_1,\cdots,\gamma_n.
	\]
	若
	\[
	\rho(f)=A,
	\qquad
	\rho(g)=B,
	\]
	即
	\[
	\bigl(f(\alpha_1)\ \cdots\ f(\alpha_l)\bigr)
	=(\beta_1\ \cdots\ \beta_m)A,
	\]
	\[
	\bigl(g(\beta_1)\ \cdots\ g(\beta_m)\bigr)
	=(\gamma_1\ \cdots\ \gamma_n)B,
	\]
	则
	\[
	\rho(gf)=BA.
	\]
\end{lemma}

\begin{proof}
	对任意 $i$，有
	\[
	\begin{aligned}
		gf(\alpha_i)
		&=g\left(\sum_{k=1}^{m}A(k,i)\beta_k\right)=\sum_{k=1}^{m}A(k,i)g(\beta_k)\\
		&=\sum_{k=1}^{m}A(k,i)\sum_{j=1}^{n}B(j,k)\gamma_j=\sum_{j=1}^{n}\left(\sum_{k=1}^{m}B(j,k)A(k,i)\right)\gamma_j\\
		&=\sum_{j=1}^{n}(BA)(j,i)\gamma_j.
	\end{aligned}
	\]
	因此 $gf$ 在上述基下的矩阵为 $BA$。
\end{proof}

\begin{corollary}
	设 $V,W$ 是有限维线性空间，$f\in\Hom(V,W)$。
	则 $f$ 是同构当且仅当 $f$ 在 $V,W$ 的基下的矩阵可逆。
\end{corollary}

\begin{proposition}
	设 $f\in\Hom(V,W)$。
	$\alpha_1,\cdots,\alpha_n$ 与 $\alpha'_1,\cdots,\alpha'_n$ 是 $V$ 的两组基，
	$\beta_1,\cdots,\beta_m$ 与 $\beta'_1,\cdots,\beta'_m$ 是 $W$ 的两组基。
	设
	\[
	(\alpha_1\ \cdots\ \alpha_n)
	=(\alpha'_1\ \cdots\ \alpha'_n)P,
	\qquad
	(\beta_1\ \cdots\ \beta_m)
	=(\beta'_1\ \cdots\ \beta'_m)Q.
	\]
	若
	\[
	\bigl(f(\alpha_1)\ \cdots\ f(\alpha_n)\bigr)
	=(\beta_1\ \cdots\ \beta_m)A,
	\]
	\[
	\bigl(f(\alpha'_1)\ \cdots\ f(\alpha'_n)\bigr)
	=(\beta'_1\ \cdots\ \beta'_m)B,
	\]
	则
	\[
	A=Q^{-1}BP.
	\]
\end{proposition}

\begin{proof}
	下证 $A$ 与 $Q^{-1}BP$ 的 $(j,i)$-元相等。
	
	\[
	\begin{aligned}
		f(\alpha_i) &= \sum_{j=1}^m A(j,i) \beta_j \\
		f(\alpha_i) &= f\left( \sum_{l=1}^n P(l,i) \alpha'_l \right) = \sum_{l=1}^n P(l,i) f(\alpha'_l)\\
			&= \sum_{l=1}^n P(l,i) \left( \sum_{t=1}^m B(t,l) \beta'_t \right) = \sum_{t=1}^m \left( \sum_{l=1}^n B(t,l) P(l,i) \right) \beta'_t \\
		&= \sum_{t=1}^m \left( \sum_{l=1}^n P(l,i) B(t,l) \right) \sum_{j=1}^m Q^{-1}(j,t) \beta_j \\
		&= \sum_{j=1}^m \left( \sum_{t,l} Q^{-1}(j,t) B(t,l) P(l,i) \right) \beta_j = \sum_{j=1}^m (Q^{-1}BP)(j,i) \beta_j
	\end{aligned}
	\]

	下面再给出一种不严格的形式化证法：
	
	\[
	\begin{aligned}
		(f(\alpha_1) \cdots f(\alpha_n)) &= (\beta_1 \cdots \beta_m) A \\
		(f(\alpha_1) \cdots f(\alpha_n)) &= f\left( (\alpha_1 \cdots \alpha_n) P \right) = (f(\alpha'_1) \cdots f(\alpha'_n)) P \\
		&= \left( (\beta'_1 \cdots \beta'_m) B \right) P = (\beta'_1 \cdots \beta'_m) (BP) \\
		&= (\beta_1 \cdots \beta_m) Q^{-1} BP
	\end{aligned}
	\]
	
	可以看到 $f$ 在 $V,W$ 下的不同的基下的矩阵相抵。
\end{proof}

\begin{proposition}
	若 $A,B\in M_{m\times n}(\F)$ 相抵，则存在
	$f\in\Hom(V,W)$，使得 $A,B$ 为 $f$ 在不同基下的矩阵。
\end{proposition}

\begin{proof}
	取定 $V$ 的基 $\alpha_1,\cdots,\alpha_n$ 与 $W$ 的基
	$\beta_1,\cdots,\beta_m$。由线性映射存在唯一性定理，存在
	$f\in\Hom(V,W)$，使
	\[
	\bigl(f(\alpha_1)\ \cdots\ f(\alpha_n)\bigr)
	=(\beta_1\ \cdots\ \beta_m)A.
	\]
	由于 $A,B$ 相抵，存在 $P\in\GL_n(\F)$、
	$Q\in\GL_m(\F)$，使得
	\[
	A=Q^{-1}BP.
	\]
	令
	\[
	(\alpha_1\ \cdots\ \alpha_n)=(\alpha'_1\ \cdots\ \alpha'_n)P,
	\qquad
	(\beta_1\ \cdots\ \beta_m)=(\beta'_1\ \cdots\ \beta'_m)Q.
	\]
	由上一个命题可知，$B$ 正是 $f$ 在基
	$\alpha'_1,\cdots,\alpha'_n$ 与 $\beta'_1,\cdots,\beta'_m$ 下的矩阵。
\end{proof}

\begin{example}
	设 $A\in M_{m\times n}(\F)$，且 $r(A)=r$。用线性映射证明：
	$A$ 相抵于
	\[
	\begin{pmatrix}
		E_r&0\\
		0&0
	\end{pmatrix}_{m\times n}.
	\]
\end{example}

\begin{proof}
	取有限维线性空间 $V,W$，以及
	$f\in\Hom(V,W)$，使得 $A$ 是 $f$ 在某两组基下的矩阵。
	此时
	\[
	\dim\ker f=n-r.
	\]
	取 $\ker f$ 的一组基 $\alpha'_{r+1},\cdots,\alpha'_n$，
	并扩充为 $V$ 的一组基
	\[
	\alpha'_1,\cdots,\alpha'_r,
	\alpha'_{r+1},\cdots,\alpha'_n.
	\]
	则 $f(\alpha'_1),\cdots,f(\alpha'_r)$ 线性无关。
	令
	\[
	\beta'_i=f(\alpha'_i),
	\qquad i=1,\cdots,r,
	\]
	并把 $\beta'_1,\cdots,\beta'_r$ 扩充为 $W$ 的一组基
	\[
	\beta'_1,\cdots,\beta'_r,
	\beta'_{r+1},\cdots,\beta'_m.
	\]
	于是
	\[
	\bigl(f(\alpha'_1)\ \cdots\ f(\alpha'_n)\bigr)
	=(\beta'_1\ \cdots\ \beta'_m)
	\begin{pmatrix}
		E_r&0\\
		0&0
	\end{pmatrix}_{m\times n}.
	\]
	故 $A$ 与该标准矩阵相抵。
\end{proof}

下面考虑 $V=W$ 且 $\dim V<\infty$ 的情形。

\begin{corollary}
	设 $f\in\End(V)$，取 $V$ 的两组基
	$\alpha_1,\cdots,\alpha_n$ 与 $\alpha'_1,\cdots,\alpha'_n$，且
	\[
	(\alpha_1\ \cdots\ \alpha_n)=(\alpha'_1\ \cdots\ \alpha'_n)P.
	\]
	若
	\[
	\bigl(f(\alpha_1)\ \cdots\ f(\alpha_n)\bigr)
	=(\alpha_1\ \cdots\ \alpha_n)A,
	\]
	\[
	\bigl(f(\alpha'_1)\ \cdots\ f(\alpha'_n)\bigr)
	=(\alpha'_1\ \cdots\ \alpha'_n)B,
	\]
	则
	\[
	A=P^{-1}BP.
	\]
\end{corollary}

\begin{definition}
	设 $f\in\Hom(V,W)$，取 $V$ 的基
	$\alpha_1,\cdots,\alpha_n$ 与 $W$ 的基 $\beta_1,\cdots,\beta_m$，
	设 $\rho(f)=A$。称 $r(A)$ 为 $f$ 的秩，记为 $r(f)$。
	显然，$r(f)$ 不依赖于基的选取。
\end{definition}

\begin{lemma}
	设
	\[
	\alpha=(\alpha_1\ \cdots\ \alpha_n)x\in V.
	\]
	则 $f(\alpha)$ 在 $\beta_1,\cdots,\beta_m$ 下的坐标为 $Ax\in\F^m$。
\end{lemma}

\begin{proof}
	若 $x=(x_1,\cdots,x_n)'$，则
	\[
	\begin{aligned}
		f(\alpha)
		&=f\left(\sum_{i=1}^{n}x_i\alpha_i\right)\\
		&=\sum_{i=1}^{n}x_i f(\alpha_i)=\sum_{i=1}^{n}x_i\left(\sum_{j=1}^{m}A(j,i)\beta_j\right)\\
		&=\sum_{j=1}^{m}\left(\sum_{i=1}^{n}A(j,i)x_i\right)\beta_j=(\beta_1\ \cdots\ \beta_m)Ax.
	\end{aligned}
	\]
\end{proof}

\begin{proposition}
	设
	\[
	\mathcal{N}(A)=\{x\in\F^n:Ax=0\},
	\qquad
	\mathcal{R}(A)=\langle \gamma_1,\cdots,\gamma_n\rangle,
	\]
	其中 $A=(\gamma_1\ \cdots\ \gamma_n)$。则
	\[
	\ker f\cong \mathcal{N}(A),
	\qquad
	\Im f\cong \mathcal{R}(A).
	\]
	从而
	\[
	\dim\ker f=n-r(A),
	\qquad
	\dim\Im f=r(A).
	\]
\end{proposition}

\begin{proof}
	对 $\alpha=(\alpha_1\ \cdots\ \alpha_n)x\in\ker f$，有
	\[
	0=f(\alpha)=(\beta_1\ \cdots\ \beta_m)Ax.
	\]
	故
	\[
	\alpha\in\ker f
	\Longleftrightarrow
	Ax=0
	\Longleftrightarrow
	x\in\mathcal{N}(A).
	\]
	于是坐标映射给出
	\[
	\ker f\cong\mathcal{N}(A).
	\]
	同理，$\Im f$ 在基 $\beta_1,\cdots,\beta_m$ 下的坐标空间
	正是 $\mathcal{R}(A)$，故
	\[
	\Im f\cong\mathcal{R}(A).
	\]
\end{proof}

于是，我们可以用新的方式对第三章中的一些结论进行重新证明。

\begin{example}
	设矩阵的尺寸使得下列乘法有意义，则有：
	\begin{enumerate}
		\item 若 $A,B\in M_{m\times n}(\F)$，则
		\[
		r(A+B)\leq r(A)+r(B).
		\]
		\item 若 $A\in M_{m\times n}(\F)$，$B\in M_{n\times s}(\F)$，则
		\[
		r(AB)\leq \min\{r(A),r(B)\}.
		\]
		\item 若 $A\in M_{m\times n}(\F)$，$B\in M_{n\times s}(\F)$，且 $AB=0$，则
		\[
		r(A)+r(B)\leq n.
		\]
		\item 若 $A\in M_{m\times n}(\F)$，$B\in M_{n\times s}(\F)$，则
		\[
		r(AB)\geq r(A)+r(B)-n.
		\]
	\end{enumerate}
\end{example}

\begin{proof}
	\begin{enumerate}
	\item 由于 $\mathcal{R}(A+B)$ 是 $\mathcal{R}(A)+\mathcal{R}(B)$ 的子空间，故
	\[
	r(A+B)=\dim\mathcal{R}(A+B)
	\leq \dim(\mathcal{R}(A)+\mathcal{R}(B))
	\leq r(A)+r(B).
	\]
	\item 设 $\rho(f)=A,\rho(g)=B$，其中
	$f\in\Hom(\F^n,\F^m)$，
	$g\in\Hom(\F^s,\F^n)$。则 $\rho(fg)=AB$。因此
	\[
	r(AB)=\dim\Im(fg)
	\leq\dim\Im f=r(A).
	\]
	又
	\[
	s-r(AB)=\dim\ker(fg)
	\geq\dim\ker g=s-r(B),
	\]
	故 $r(AB)\leq \min\{r(A),r(B)\}$。
	\item 记号同上。若 $AB=0$，则
	\[
	fg=0
	\Longleftrightarrow
	\Im g\subset\ker f.
	\]
	从而
	\[
	r(B)=\dim\Im g
	\leq \dim\ker f=n-r(A),
	\]
	即 $r(A)+r(B)\leq n$。
	\item 因 $\mathcal{N}(B)\subset\mathcal{N}(AB)$，取
	$\mathcal{N}(B)$ 的基 $\varepsilon_1,\cdots,\varepsilon_b$，
	并扩充为 $\mathcal{N}(AB)$ 的基
	\[
	\varepsilon_1,\cdots,\varepsilon_b,
	\varepsilon_{b+1},\cdots,\varepsilon_c.
	\]
	则
	\[
	B\varepsilon_{b+1},\cdots,B\varepsilon_c\in\mathcal{N}(A)
	\]
	且它们线性无关。故
	\[
	c-b\leq \dim\mathcal{N}(A)=n-r(A).
	\]
	又
	\[
	c=s-r(AB),
	\qquad
	b=s-r(B),
	\]
	所以
	\[
	r(AB)\geq r(A)+r(B)-n.
	\]
\end{enumerate}
\end{proof}

\section{线性空间的子空间分解}

\begin{example}
	设 $T\in\End(V)$，且
	\[
	T^2-3T+2\id_V=0.
	\]
	记
	\[
	V_1=\ker(T-2\id_V),
	\qquad
	V_2=\ker(T-\id_V).
	\]
	证明
	\[
	V=V_1\oplus V_2.
	\]
\end{example}

\begin{proof}
	先证 $V=V_1+V_2$。对任意 $\alpha\in V$，有
	\[
	\alpha=(T\alpha-\alpha)-(T\alpha-2\alpha).
	\]
	并且
	\[
	(T-2\id_V)(T\alpha-\alpha)
	=(T^2-3T+2\id_V)\alpha=0,
	\]
	故 $T\alpha-\alpha\in V_1$；又
	\[
	(T-\id_V)(T\alpha-2\alpha)
	=(T^2-3T+2\id_V)\alpha=0,
	\]
	故 $T\alpha-2\alpha\in V_2$。从而 $V=V_1+V_2$。
	
	再证和为直和。设
	\[
	0=\alpha_1+\alpha_2,
	\qquad
	\alpha_1\in V_1,
	\quad
	\alpha_2\in V_2.
	\]
	则 $-\alpha_1=\alpha_2$。两边作用 $T$，得
	\[
	-T\alpha_1=T\alpha_2,
	\qquad
	-2\alpha_1=\alpha_2.
	\]
	又由 $\alpha_2=-\alpha_1$，得 $\alpha_1=0$，进而 $\alpha_2=0$。
	故 $V_1+V_2$ 是直和。
\end{proof}

因此，我们从零化多项式的角度考虑子空间分解。

\begin{lemma}
	设 $\dim V=n<\infty$，$T\in\End(V)$，则 $T$ 存在零化多项式。
\end{lemma}

\begin{proof}
	由于
	\[
	\dim\End(V)=n^2,
	\]
	故
	\[
	\id_V,T,T^2,\cdots,T^{n^2}
	\]
	线性相关。于是存在非零多项式 $h(x)\in\F[x]$，使得$ h(T)=0.$
\end{proof}

\begin{definition}
	设 $m\in\F[x]$，若 $m(T)=0$，且 $m$ 是所有满足该条件的非零多项式中次数最低的首一多项式，则称 $m$ 为 $T$ 的最小多项式，记为 $m_T$。
\end{definition}

\begin{lemma}
	最小多项式存在且唯一。
	\[\{T\text{的零化多项式}\}=\{mh:h\in \F[x],m = m_T\} \]
\end{lemma}

\begin{proof}
	存在性由上一引理给出。若 $p,q$ 都是 $T$ 的最小多项式，则 $p,q$ 均首一且次数相同。若 $p\neq q$，则
	\[
	(p-q)(T)=0,
	\]
	而 $p-q$ 是次数更低的非零零化多项式，矛盾。故 $p=q$。
\end{proof}

\begin{lemma}
	设 $T\in\End(V)$，$T$ 在某组基下的矩阵为 $A$，则 $T$ 与 $A$ 有相同的最小多项式。
\end{lemma}

\begin{proof}
	设 $\rho(T)=A$，则对任意多项式 $g(x)=\sum_{i=0}^{k}a_i x^i\in\F[x]$，有
	\[
	\rho(g(T))=\sum_{i=0}^{k}a_i A^i=g(A).
	\]
	由于 $\rho$ 是同构，$g(T)=0$ 当且仅当 $g(A)=0$。
	因此 $T$ 与 $A$ 有完全相同的零化多项式，从而有相同的最小多项式。
\end{proof}

\begin{theorem}(准素分解定理)
	设 $V$ 是 $\F$ 上的线性空间，$T\in\End(V)$。设
	\[
	g=g_1\cdots g_s,
	\]
	其中 $g_1,\cdots,g_s\in\F[x]$ 两两互素，则
	\[
	\ker g(T)=\ker g_1(T)\oplus\cdots\oplus\ker g_s(T).
	\]
	若 $g(T)=0$，则
	\[
	V=\ker g_1(T)\oplus\cdots\oplus\ker g_s(T).
	\]
\end{theorem}

\begin{proof}
	记
	\[
	h_i=\frac{g}{g_i},
	\qquad i=1,\cdots,s.
	\]
	由于 $g_1,\cdots,g_s$ 两两互素，故 $h_1,\cdots,h_s$ 互素，存在 $u_1,\cdots,u_s\in\F[x]$，使得
	\[
	u_1h_1+\cdots+u_sh_s=1.
	\]
	
	先证和空间等于 $\ker g(T)$。若 $\alpha\in\ker g_i(T)$，则
	\[
	g(T)\alpha=h_i(T)g_i(T)\alpha=0,
	\]
	故
	\[
	\ker g_1(T)+\cdots+\ker g_s(T)\subset \ker g(T).
	\]
	反过来，若 $\alpha\in\ker g(T)$，则
	\[
	\alpha=u_1(T)h_1(T)\alpha+\cdots+u_s(T)h_s(T)\alpha.
	\]
	并且
	\[
	g_i(T)u_i(T)h_i(T)\alpha=u_i(T)g(T)\alpha=0,
	\]
	故 $u_i(T)h_i(T)\alpha\in\ker g_i(T)$。于是
	\[
	\ker g(T)=\ker g_1(T)+\cdots+\ker g_s(T).
	\]
	
	再证直和。设
	\[
	\beta_1+\cdots+\beta_s=0,
	\qquad
	\beta_i\in\ker g_i(T).
	\]
	固定 $i$。由于 $(h_i,g_i)=1$，存在 $p_i,q_i\in\F[x]$，使得
	\[
	p_ih_i+q_ig_i=1.
	\]
	又 $h_i(T)\beta_j=0$ 对 $j\neq i$ 成立，所以
	\[
	h_i(T)\beta_i
	=-\sum_{j\neq i}h_i(T)\beta_j=0.
	\]
	于是
	\[
	\beta_i=\bigl(p_i(T)h_i(T)+q_i(T)g_i(T)\bigr)\beta_i=0.
	\]
	故每个 $\beta_i=0$，从而上述和为直和。
	
	若 $g(T)=0$，则 $\ker g(T)=V$，于是得到
	\[
	V=\ker g_1(T)\oplus\cdots\oplus\ker g_s(T).
	\]
\end{proof}
\newpage
\section{一些有价值的习题}

\begin{problem}(线性相关与线性无关)
	设 $S_1,S_2$ 是线性无关集。
	\begin{enumerate}
		\item $S_1\cap S_2$ 是否线性无关？
		\item $S_1\cup S_2$ 是否一定线性无关？
		\item 用 $\langle S_1\rangle\cap\langle S_2\rangle$ 刻画 $S_1\cup S_2$ 何时线性无关。
	\end{enumerate}
\end{problem}

\begin{proof}
	\begin{enumerate}
		\item $S_1\cap S_2$ 是 $S_1$ 的子集，故若 $S_1\cap S_2$ 非空，则它线性无关。
		\item 不一定。例如在 $\R^2$ 中，取
		\[
		S_1=\left\{
		\begin{pmatrix}1\\0\end{pmatrix},
		\begin{pmatrix}0\\1\end{pmatrix}
		\right\},
		\qquad
		S_2=\left\{
		\begin{pmatrix}1\\1\end{pmatrix}
		\right\}.
		\]
		则 $S_1,S_2$ 都线性无关，但 $S_1\cup S_2$ 线性相关。
		\item 设
		\[
		S_1=\{\alpha_1,\cdots,\alpha_m\},
		\qquad
		S_2=\{\beta_1,\cdots,\beta_n\},
		\]
		且 $S_1,S_2$ 中列出的向量互不相同。则
		\[
		S_1\cup S_2\text{ 线性无关}
		\Longleftrightarrow
		\langle S_1\rangle\cap\langle S_2\rangle=\{0\}.
		\]
		必要性：设 $\alpha\in\langle S_1\rangle\cap\langle S_2\rangle$，则存在 $x_i,y_j\in\F$，使
		\[
		\alpha=\sum_{i=1}^{m}x_i\alpha_i=\sum_{j=1}^{n}y_j\beta_j.
		\]
		于是
		\[
		\sum_{i=1}^{m}x_i\alpha_i-\sum_{j=1}^{n}y_j\beta_j=0.
		\]
		由 $S_1\cup S_2$ 线性无关，得 $x_i=y_j=0$，故 $\alpha=0$。
		
		充分性：设
		\[
		a_1\alpha_1+\cdots+a_m\alpha_m+b_1\beta_1+\cdots+b_n\beta_n=0.
		\]
		则
		\[
		a_1\alpha_1+\cdots+a_m\alpha_m
		=-(b_1\beta_1+\cdots+b_n\beta_n)
		\in \langle S_1\rangle\cap\langle S_2\rangle=\{0\}.
		\]
		由 $S_1,S_2$ 分别线性无关，得 $a_i=0,b_j=0$。故 $S_1\cup S_2$ 线性无关。
	\end{enumerate}
\end{proof}

\begin{problem}
	令
	\[
	\Q(\sqrt[n]{2})=\{c=f(\sqrt[n]{2}):f(x)\in\Q[x]\}.
	\]
	证明：
	\begin{enumerate}
		\item $\Q(\sqrt[n]{2})$ 是 $\Q$ 上的线性空间。事实上，$\Q(\sqrt[n]{2})$ 是一个数域。
		\item 有
		\[
		\Q(\sqrt[n]{2})=
		\left\{
		\sum_{i=0}^{n-1}a_i(\sqrt[n]{2})^i:a_i\in\Q
		\right\}.
		\]
		\item $\dim_{\Q}\Q(\sqrt[n]{2})=n$，并求出一组基。
	\end{enumerate}
\end{problem}

\begin{proof}
	记 $r=\sqrt[n]{2}$。
	\begin{enumerate}
		\item 首先 $0=0(r)\in\Q(r)$。若 $c=f(r),d=g(r)\in\Q(r)$，则
		\[
		c+d=(f+g)(r)\in\Q(r),
		\qquad
		cd=(fg)(r)\in\Q(r).
		\]
		若 $k\in\Q$，则
		\[
		kc=(kf)(r)\in\Q(r).
		\]
		故 $\Q(r)$ 对加法与数乘封闭，从而是 $\Q$ 上的线性空间。
		
		进一步，若 $0\neq c=f(r)$，由下面的第二部分可设 $\deg f<n$。多项式 $x^n-2$ 由 Eisenstein 判别法在 $\Q$ 上不可约。若 $f\neq 0$，则 $(f,x^n-2)=1$，故存在 $u,v\in\Q[x]$，使
		\[
		uf+v(x^n-2)=1.
		\]
		代入 $r$ 得
		\[
		u(r)c=1,
		\]
		故 $c^{-1}=u(r)\in\Q(r)$。因此 $\Q(r)$ 是一个数域。
		
		\item 对任意 $c=f(r)\in\Q(r)$，对 $f(x)$ 关于 $x^n-2$ 作带余除法：
		\[
		f(x)=q(x)(x^n-2)+s(x),
		\qquad s=0\text{ 或 }\deg s<n.
		\]
		代入 $r$ 得
		\[
		c=f(r)=s(r).
		\]
		于是 $c$ 可写成
		\[
		c=\sum_{i=0}^{n-1}a_ir^i,
		\qquad a_i\in\Q.
		\]
		反向包含显然成立，故等式成立。
		
		\item 由第二部分，$1,r,r^2,\cdots,r^{n-1}$ 生成 $\Q(r)$。若
		\[
		\sum_{i=0}^{n-1}a_ir^i=0,
		\qquad a_i\in\Q,
		\]
		令 $h(x)=\sum_{i=0}^{n-1}a_ix^i$。若 $h\neq0$，则 $h$ 与 $x^n-2$ 有公共根 $r$。由于 $x^n-2$ 在 $\Q$ 上不可约，由前面关于公共根的命题可知 $x^n-2\mid h$，这与 $\deg h<n$ 矛盾。故 $h=0$，即 $a_i=0$。
		
		因此
		\[
		1,r,r^2,\cdots,r^{n-1}
		\]
		是 $\Q(r)$ 的一组基，从而
		\[
		\dim_{\Q}\Q(\sqrt[n]{2})=n.
		\]
	\end{enumerate}
\end{proof}

\begin{definition}
	通常线性代数中所说的基是 Hamel 基。设 $B\subset V$，若满足：
	\begin{enumerate}
		\item $B$ 中任意有限多个不同向量线性无关；
		\item 任意 $v\in V$ 都能由 $B$ 中有限多个向量线性表示，
	\end{enumerate}
	则 $B$ 是 $V$ 的一组 Hamel 基。也就是说，在线性代数意义下，只有有限多个系数可以不为 $0$。
\end{definition}

	
	例如在多项式空间 $P(\R)$ 中，
	\[
	1,x,x^2,x^3,\cdots
	\]
	是一组 Hamel 基。多项式
	\[
	p(x)=3+2x-5x^4
	\]
	的坐标为
	\[
	(3,2,0,0,-5,0,\cdots),
	\]
	其中只有有限多项不为 $0$。
	\begin{definition}
			Schauder 基是泛函分析中的概念，它允许向量表示成无穷级数，但要求级数在给定范数下收敛。即
		\[
		x=\sum_{n=1}^{\infty}a_ne_n,
		\qquad
		\lim_{N\to\infty}\left\|x-\sum_{n=1}^{N}a_ne_n\right\|=0.
		\]
	\end{definition}

	因此，几何级数
	\[
	1+x+x^2+\cdots=\frac{1}{1-x},
	\qquad x\in(-1,1),
	\]
	并不给出集合
	\[
	S=\left\{g(x)=\frac{1}{1-x},\ f_0(x)=1,\ f_1(x)=x,\ f_2(x)=x^2,\cdots\right\}
	\]
	在 Hamel 意义下的线性相关关系，因为这不是有限线性组合；在 Schauder 意义下才可以讨论这种无穷展开。


\begin{problem}[线段与凸集]
	设 $u,v\in\R^n$，连接 $u,v$ 的线段定义为
	\[
	\ell=\{tu+(1-t)v:t\in[0,1]\}.
	\]
	\begin{enumerate}
		\item 若 $h:\R^n\to\R^m$ 是线性映射，证明 $h(\ell)$ 是连接 $h(u),h(v)$ 的线段。
		\item 一个$\R^n$的子集是凸的，若对其中任意两个点，两点连成的线段是这个子集的子集。若 $S\subset\R^n$ 是凸集，证明线性映射把凸集映为凸集。
	\end{enumerate}
\end{problem}

\begin{proof}
	\begin{enumerate}
		\item 对任意 $\alpha=tu+(1-t)v\in\ell$，由于 $h$ 是线性映射，所以
		\[
		h(\alpha)=t h(u)+(1-t)h(v).
		\]
		因此 $h(\alpha)$ 位于以 $h(u),h(v)$ 为端点的线段上。
		\item 设 $f:\R^n\to\R^m$ 是线性映射，$S\subset\R^n$ 是凸集。任取 $\alpha,\beta\in f(S)$，则存在 $\alpha',\beta'\in S$，使得
		\[
		\alpha=f(\alpha'),
		\qquad
		\beta=f(\beta').
		\]
		对任意 $t\in[0,1]$，由 $S$ 的凸性有
		\[
		t\alpha'+(1-t)\beta'\in S.
		\]
		从而
		\[
		t\alpha+(1-t)\beta
		=t f(\alpha')+(1-t)f(\beta')
		=f(t\alpha'+(1-t)\beta')\in f(S).
		\]
		故 $f(S)$ 也是凸集。
	\end{enumerate}
\end{proof}

\begin{problem}
	设 $V_1,V_2,V_3$ 是线性空间 $L$ 的子空间，且 $V_1\subset V_3$。证明
	\[
	V_1+(V_2\cap V_3)=(V_1+V_2)\cap V_3.
	\]
\end{problem}

\begin{proof}
	先证
	\[
	V_1+(V_2\cap V_3)\subset (V_1+V_2)\cap V_3.
	\]
	任取 $x\in V_1+(V_2\cap V_3)$，则
	\[
	x=x_1+x_2,
	\qquad x_1\in V_1,
	\quad x_2\in V_2\cap V_3.
	\]
	因为 $V_1\subset V_3$，所以 $x_1\in V_3$，从而 $x\in V_3$，且 $x\in V_1+V_2$。故
	$x\in (V_1+V_2)\cap V_3$。
	
	反过来，任取 $x\in (V_1+V_2)\cap V_3$，则
	\[
	x=x_1+x_2,
	\qquad x_1\in V_1,
	\quad x_2\in V_2,
	\]
	且 $x\in V_3$。由于 $V_1\subset V_3$，故 $x_1\in V_3$，于是
	\[
	x_2=x-x_1\in V_3.
	\]
	所以 $x_2\in V_2\cap V_3$，从而 $x\in V_1+(V_2\cap V_3)$。
	
	综上，结论成立。
\end{proof}

\begin{problem}
	求 $V_1+V_2$ 与 $V_1\cap V_2$ 的一组基，并计算维数。
	\begin{enumerate}
		\item $V_1=L(\alpha_1,\alpha_2),V_2=L(\beta_1,\beta_2)$，其中
		\[
		\alpha_1=(1,3,2,4),\quad
		\alpha_2=(2,1,0,1),
		\]
		\[
		\beta_1=(3,2,1,0),\quad
		\beta_2=(-1,2,2,3).
		\]
		\item $V_1=L(\alpha_1,\alpha_2,\alpha_3),V_2=L(\beta_1,\beta_2)$，其中
		\[
		\alpha_1=(0,2,0,-3),
		\quad
		\alpha_2=(2,1,2,0),
		\quad
		\alpha_3=(1,4,-1,-5),
		\]
		\[
		\beta_1=(-2,-5,6,5),
		\quad
		\beta_2=(3,3,1,-8).
		\]
	\end{enumerate}
\end{problem}

\begin{proof}
	\begin{enumerate}
		\item 对矩阵 $A=(\alpha_1\ \alpha_2\ \beta_1\ \beta_2)$ 作初等行变换：
		\[
		A=
		\begin{pmatrix}
			1&2&3&-1\\
			3&1&2&2\\
			2&0&1&2\\
			4&1&0&3
		\end{pmatrix}
		\sim
		\begin{pmatrix}
			1&0&0&1\\
			0&1&0&-1\\
			0&0&1&0\\
			0&0&0&0
		\end{pmatrix}.
		\]
		由此可知主元列对应 $\alpha_1,\alpha_2,\beta_1$，故
		\[
		\alpha_1,\alpha_2,\beta_1
		\]
		是 $V_1+V_2$ 的一组基，
		\[
		\dim(V_1+V_2)=3.
		\]
		又 $\dim V_1=2,\dim V_2=2$，由维数公式
		\[
		\dim(V_1\cap V_2)=2+2-3=1.
		\]
		由行变换得到关系
		\[
		-\alpha_1+\alpha_2+\beta_2=0,
		\]
		即
		\[
		\beta_2=\alpha_1-\alpha_2\in V_1.
		\]
		同时 $\beta_2\in V_2$，且 $\beta_2\neq0$，故
		\[
		V_1\cap V_2=L(\beta_2),
		\qquad
		\dim(V_1\cap V_2)=1.
		\]
		
		\item 对矩阵 $A=(\alpha_1\ \alpha_2\ \alpha_3\ \beta_1\ \beta_2)$ 作初等行变换：
		\[
		A=
		\begin{pmatrix}
			0&2&1&-2&3\\
			2&1&4&-5&3\\
			0&2&-1&6&1\\
			-3&0&-5&5&-8
		\end{pmatrix}
		\sim
		\begin{pmatrix}
			1&0&0&5&0\\
			0&1&0&1&0\\
			0&0&1&-4&0\\
			0&0&0&0&1
		\end{pmatrix}.
		\]
		所以主元列对应 $\alpha_1,\alpha_2,\alpha_3,\beta_2$，故
		\[
		\alpha_1,\alpha_2,\alpha_3,\beta_2
		\]
		是 $V_1+V_2$ 的一组基，
		\[
		\dim(V_1+V_2)=4.
		\]
		又 $\dim V_1=3,\dim V_2=2$，所以
		\[
		\dim(V_1\cap V_2)=3+2-4=1.
		\]
		由行变换可得
		\[
		-5\alpha_1-\alpha_2+4\alpha_3+\beta_1=0,
		\]
		即
		\[
		\beta_1=5\alpha_1+\alpha_2-4\alpha_3\in V_1.
		\]
		而 $\beta_1\in V_2$ 且 $\beta_1\neq0$，故
		\[
		V_1\cap V_2=L(\beta_1),
		\qquad
		\dim(V_1\cap V_2)=1.
		\]
	\end{enumerate}
\end{proof}

\begin{algo}[求子空间的和与交]
	设 $V_1,V_2$ 是 $L$ 的子空间，取定 $L$ 的一组基
	\[
	\xi_1,\cdots,\xi_n.
	\]
	设 $V_1$ 由 $\alpha_1,\cdots,\alpha_p$ 生成，$V_2$ 由 $\beta_1,\cdots,\beta_q$ 生成。若
	\[
	(\alpha_1\ \cdots\ \alpha_p)=(\xi_1\ \cdots\ \xi_n)A,
	\qquad
	(\beta_1\ \cdots\ \beta_q)=(\xi_1\ \cdots\ \xi_n)B,
	\]
	则求 $V_1+V_2$ 的一组基，等价于求矩阵
	\[
	C=(A\ B)
	\]
	的列向量组的极大无关组。具体地，对 $C$ 作初等行变换，阶梯形矩阵中的主元列所对应的原向量，就是 $V_1+V_2$ 的一组基。
	
	求 $V_1\cap V_2$ 时，可以解方程组
	\[
	Ax=By,
	\qquad \text{即}\qquad
	(A\ -B)
	\begin{pmatrix}
		x\\y
	\end{pmatrix}=0.
	\]
	每一个解都给出一个交空间中的向量
	\[
	(\xi_1\ \cdots\ \xi_n)Ax=(\xi_1\ \cdots\ \xi_n)By.
	\]
	取这些向量的极大无关组，就得到 $V_1\cap V_2$ 的一组基。维数也可以用
	\[
	\dim(V_1\cap V_2)=\dim V_1+\dim V_2-\dim(V_1+V_2)
	\]
	计算。
\end{algo}

\begin{theorem}
	设 $V$ 是数域 $\F$ 上的非零线性空间，且 $\F$ 是无限域。若 $V_1,\cdots,V_s$ 是 $V$ 的有限多个真子空间，则
	\[
	V\neq\bigcup_{i=1}^{s}V_i.
	\]
	换言之，$V$ 不可能等于有限多个真子空间的并。
\end{theorem}

\begin{proof}
	对真子空间的个数 $s$ 作归纳。
	
	当 $s=2$ 时，若 $V=V_1\cup V_2$，取 $\alpha\in V_1\setminus V_2$，$\beta\in V_2\setminus V_1$。则 $\alpha+\beta$ 既不属于 $V_1$，也不属于 $V_2$，矛盾。
	
	假设结论对任意 $s$ 个真子空间成立。设 $V_1,\cdots,V_s,V_{s+1}$ 是 $V$ 的 $s+1$ 个真子空间。由归纳假设，存在
	\[
	\alpha\in V,
	\qquad
	\alpha\notin\bigcup_{j=1}^{s}V_j.
	\]
	若 $\alpha\notin V_{s+1}$，结论已经成立。否则 $\alpha\in V_{s+1}$。由于 $V_{s+1}$ 是真子空间，可取
	\[
	\beta\notin V_{s+1}.
	\]
	考虑集合
	\[
	W=\{\lambda\alpha+\beta:\lambda\in\F\}.
	\]
	若 $\lambda_1\alpha+\beta=\lambda_2\alpha+\beta$，则 $(\lambda_1-\lambda_2)\alpha=0$，从而 $\lambda_1=\lambda_2$。因为 $\F$ 有无穷多个元素，故 $W$ 有无穷多个元素。
	
	对每个 $1\leq j\leq s$，$W\cap V_j$ 至多含有一个元素。否则若
	\[
	\lambda_1\alpha+\beta,
	\quad
	\lambda_2\alpha+\beta\in W\cap V_j,
	\quad
	\lambda_1\neq\lambda_2,
	\]
	则
	\[
	(\lambda_1-\lambda_2)\alpha\in V_j,
	\]
	从而 $\alpha\in V_j$，矛盾。因此 $\bigcup_{j=1}^{s}V_j$ 至多包含 $W$ 中有限多个元素。
	
	又 $W\cap V_{s+1}=\varnothing$。事实上，若 $\lambda\alpha+\beta\in V_{s+1}$，由于 $\lambda\alpha\in V_{s+1}$，可得 $\beta\in V_{s+1}$，矛盾。
	
	所以可以取 $\gamma\in W$，使
	\[
	\gamma\notin\bigcup_{j=1}^{s+1}V_j.
	\]
	归纳完成。
\end{proof}

\begin{problem}
	固定实数 $x_0$，令
	\[
	V=\{f(x)\in\R_n[x]:f(x_0)=0\}.
	\]
	\begin{enumerate}
		\item 证明 $V$ 是 $\R_n[x]$ 的子空间。
		\item 求 $V$ 在 $\R_n[x]$ 中的一个直和补 $W$，即求 $W$ 使
		\[
		\R_n[x]=V\oplus W,
		\]
		并写出投影 $p_W$。
		\item 若 $x_0\in[a,b]$，把上述结论推广到 $C[a,b]$。
	\end{enumerate}
\end{problem}

\begin{proof}
	\begin{enumerate}
		\item 首先 $0\in V$。若 $f,g\in V$，$a,b\in\R$，则
		\[
		(af+bg)(x_0)=af(x_0)+bg(x_0)=0.
		\]
		同时 $af+bg\in\R_n[x]$，故 $af+bg\in V$。所以 $V$ 是 $\R_n[x]$ 的子空间。
		
		事实上，$V$ 中的元素正是以 $x_0$ 为根的多项式，因此可以写为
		\[
		(x-x_0)g(x)
		\]
		的形式。故
		\[
		x-x_0,(x-x_0)^2,\cdots,(x-x_0)^{n-1}
		\]
		是 $V$ 的一组基，从而 $\dim V=n-1$。
		
		\item 对任意 $f(x)\in\R_n[x]$，作带余除法可得
		\[
		f(x)=(x-x_0)q(x)+f(x_0).
		\]
		其中 $(x-x_0)q(x)\in V$，$f(x_0)$ 是常数多项式。于是
		\[
		\R_n[x]=V+\R.
		\]
		若 $f\in V\cap\R$，则 $f$ 是常数多项式且 $f(x_0)=0$，所以 $f=0$。故
		\[
		\R_n[x]=V\oplus\R.
		\]
		因此可取
		\[
		W=\R.
		\]
		此时到 $W$ 上的投影为
		\[
		p_W(f)=f(x_0)
		\]
		这里右端视为常数多项式。
		
		\item 在 $C[a,b]$ 中同理，令
		\[
		V=\{f(x)\in C[a,b]:f(x_0)=0\}.
		\]
		对任意 $f\in C[a,b]$，有分解
		\[
		f(x)=\bigl(f(x)-f(x_0)\bigr)+f(x_0),
		\]
		其中 $f(x)-f(x_0)\in V$，$f(x_0)\in\R$。若 $g\in V\cap\R$，则 $g$ 是常数函数且 $g(x_0)=0$，故 $g=0$。因此
		\[
		C[a,b]=V\oplus\R.
		\]
	\end{enumerate}
\end{proof}
\begin{proposition}
		一般地，设 $V_1,\cdots,V_s$ 是 $L$ 的子空间，则以下条件等价：
	\begin{enumerate}
		\item $V_1+\cdots+V_s$ 是直和；
		\item 若
		\[
		0=\alpha_1+\cdots+\alpha_s,
		\qquad \alpha_i\in V_i,
		\]
		则 $\alpha_i=0$，$1\leq i\leq s$，即 $0$ 的分解唯一；
		\item 对任意 $i$，有
		\[
		V_i\cap\left(\sum_{1\leq j\leq s,\ j\neq i}V_j\right)=0;
		\]
		\item 把每个 $V_i$ 的基合起来，可以得到 $V_1+\cdots+V_s$ 的一组基。
	\end{enumerate}
	进一步，如果每个 $V_i$ 都是有限维的，则
	\[
	V_1+\cdots+V_s\text{ 是直和}
	\Longleftrightarrow
	\dim(V_1+
	\cdots+V_s)=\dim V_1+
	\cdots+
	\dim V_s.
	\]
\end{proposition}


\begin{problem}[陪集与商空间]
	设 $U$ 是线性空间 $V$ 的子空间。对 $v\in V$，定义 $U$ 在 $V$ 中含有 $v$ 的陪集为
	\[
	v+U=\{v+u:u\in U\}.
	\]
	也记作 $\overline v$ 或 $v\pmod U$。
	\begin{enumerate}
		\item 设 $V=\F^n$，$A\in M_{m\times n}(\F)$。取 $U$ 为齐次线性方程组 $AX=0$ 的解空间。若 $AX=\beta$ 有解，证明它的解集是 $V$ 中 $U$ 的一个陪集。
		\item 设 $f:V\to W$ 是线性映射。若 $\beta\in\Im f$，证明
		\[
		f^{-1}(\beta)=\{\alpha\in V:f(\alpha)=\beta\}
		\]
		是 $\ker f$ 的一个陪集。
		\item 证明
		\[
		\overline u=\overline v
		\Longleftrightarrow
		u-v\in U.
		\]
		\item 记 $V/U$ 为 $U$ 在 $V$ 中所有陪集的集合。定义
		\[
		\overline u+
		\overline v=\overline{u+v},
		\qquad
		k\overline u=\overline{ku}.
		\]
		证明这些运算良定，并证明 $V/U$ 是 $\F$ 上的线性空间。
		\item 若 $W$ 是 $U$ 的直和补，即 $V=U\oplus W$，证明
		\[
		W\cong V/U.
		\]
		\item 若 $\dim V<\infty$，证明
		\[
		\dim(V/U)=\dim V-\dim U.
		\]
	\end{enumerate}
\end{problem}

\begin{proof}
	\begin{enumerate}
		\item 设 $\gamma$ 是 $AX=\beta$ 的一个解。由线性方程组解的结构定理，$AX=\beta$ 的解集为
		\[
		S=\{\gamma+\xi:\xi\in U\}=\gamma+U.
		\]
		所以它是 $U$ 的一个陪集。
		
		\item 任取 $\gamma\in f^{-1}(\beta)$，只需证明
		\[
		f^{-1}(\beta)=\gamma+
		\ker f.
		\]
		若 $\eta\in\gamma+
		\ker f$，则
		\[
		f(\eta)=f(\gamma)+0=\beta,
		\]
		故 $\eta\in f^{-1}(\beta)$。反过来，若 $\eta\in f^{-1}(\beta)$，则
		\[
		f(\eta-
		\gamma)=f(\eta)-f(\gamma)=\beta-
		\beta=0,
		\]
		故 $\eta-
		\gamma\in\ker f$，于是 $\eta=\gamma+(\eta-
		\gamma)\in\gamma+
		\ker f$。
		
		\item 由定义，$\overline u=\overline v$ 等价于 $u+U=v+U$。这又等价于 $u\in v+U$，即存在 $\xi\in U$，使 $u=v+\xi$。于是
		\[
		u-v=\xi\in U.
		\]
		反向同理。因此
		\[
		\overline u=\overline v
		\Longleftrightarrow
		u-v\in U.
		\]
		
		\item 先证良定。若 $\overline u=\overline{u'}$，$\overline v=\overline{v'}$，则
		\[
		u-u'\in U,
		\qquad
		v-v'\in U.
		\]
		因此
		\[
		(u+v)-(u'+v')=(u-u')+(v-v')\in U,
		\]
		所以
		\[
		\overline{u+v}=\overline{u'+v'}.
		\]
		类似地，若 $\overline u=\overline{u'}$，则
		\[
		ku-ku'=k(u-u')\in U,
		\]
		所以 $\overline{ku}=\overline{ku'}$。故加法与数乘均良定。
		
		下面验证线性空间公理。对任意 $u,v,w\in V$，$k,l\in\F$，有
		\[
		\overline u+
		\overline v=\overline{u+v}=\overline{v+u}=\overline v+
		\overline u,
		\]
		\[
		(\overline u+
		\overline v)+\overline w=
		\overline{(u+v)+w}=\overline{u+(v+w)}=
		\overline u+(\overline v+
		\overline w),
		\]
		\[
		\overline 0+
		\overline u=\overline u,
		\qquad
		\overline u+\overline{-u}=\overline0,
		\]
		\[
		1\overline u=\overline u,
		\qquad
		k(l\overline u)=(kl)\overline u,
		\]
		\[
		k(\overline u+
		\overline v)=k\overline u+k\overline v,
		\qquad
		(k+l)\overline u=k\overline u+l\overline u.
		\]
		所以 $V/U$ 是一个 $\F$ 上的线性空间，称为 $V$ 关于 $U$ 的商空间。
		
		\item 定义映射
		\[
		\varphi:W\to V/U,
		\qquad
		w\longmapsto \overline w=w+U.
		\]
		显然 $\varphi$ 线性。若 $\varphi(w)=\overline0$，则 $w\in U$。又 $w\in W$，且 $V=U\oplus W$，故 $w=0$。所以 $\varphi$ 单射。
		
		任取 $\overline v\in V/U$。由于 $V=U\oplus W$，存在 $u\in U,w\in W$，使
		\[
		v=u+w.
		\]
		由第三部分，$v-w=u\in U$，故
		\[
		\overline v=\overline w=\varphi(w).
		\]
		所以 $\varphi$ 满射。故 $W\cong V/U$。
		
		\item 因 $\dim V<\infty$，可取 $U$ 在 $V$ 中的一个补空间 $W$，使
		\[
		V=U\oplus W.
		\]
		由上一部分有 $W\cong V/U$，因此
		\[
		\dim(V/U)=\dim W=\dim V-
		\dim U.
		\]
	\end{enumerate}
\end{proof}
